\chapter{Spaces of Constant Curvature}
\label{chap:dc-spaces-of-constant-curvature}

Faithful transcription of do Carmo, \emph{Riemannian Geometry}. Each numbered
statement keeps do Carmo's number, kind and (name); proofs keep his explicit
cross-references ("by \cref{thm:dc-ch8-2-1}", "from \cref{cor:dc-ch5-2-5}", "Lemma
3.3 of Chap. 7") so the dependency graph is recoverable from the text.

\section{Introduction}

Among the Riemannian manifolds, those with constant sectional curvature are the
most simple. They are related to classical non-Euclidean geometry, which is
historically the first example of a geometric structure different from the
Euclidean. An important property of the spaces with constant curvature is that
they have a sufficiently large number of local isometries (cf. \cref{cor:dc-ch8-2-2});
this property of "free mobility of small triangles" was considered as a postulate
to be satisfied by the non-Euclidean geometries.

When we multiply a Riemannian metric by a positive constant $c$, then its
sectional curvature is multiplied by $1/c$. Therefore, up to a similarity, we can
suppose that the value of the constant sectional curvature of a Riemannian
manifold is $1$, $0$, or $-1$; this hypothesis is assumed without further comment.
So far we have encountered two examples of Riemannian manifolds with constant
sectional curvature $K$: the Euclidean space $\mathbb{R}^n$ with $K\equiv 0$ and
the unit sphere $S^n\subset\mathbb{R}^{n+1}$ with $K\equiv 1$. In this chapter we
introduce a Riemannian manifold, the \emph{hyperbolic space} $H^n$ of dimension $n$,
which has sectional curvature $K\equiv -1$. The manifolds $\mathbb{R}^n$, $S^n$
and $H^n$ are complete and simply connected (cf. Section 3). The main theorem of
this chapter is that these are essentially the only complete, simply connected
Riemannian manifolds with constant sectional curvature. This allows us to reduce
the problem of finding all the complete manifolds of constant curvature to the
problem of determining certain subgroups of the group of isometries of
$\mathbb{R}^n$, $S^n$ and $H^n$ (cf. Section 4).

To prove this, we introduce in Section 2 a slightly more general theorem on the
determination of the metric by means of the curvature. We frequently use the
classical expression \emph{spaces of constant curvature} to designate the Riemannian
manifolds of constant sectional curvature. In the last section we describe the
isometries of the hyperbolic space $H^n$ and identify certain important
hypersurfaces of $H^n$, namely, the horospheres and the hyperspheres.

\section{Theorem of Cartan on the determination of the metric by means of the curvature}

Let $M$ and $\tilde{M}$ be two Riemannian manifolds of dimension $n$ and let
$p\in M$ and $\tilde{p}\in\tilde{M}$. Choose a linear isometry
$i:T_p(M)\to T_{\tilde{p}}(\tilde{M})$. Let $V\subset M$ be a normal neighborhood
of $p$ such that $\exp_{\tilde{p}}$ is defined at $i\circ\exp_p^{-1}(V)$. Define a
mapping $f:V\to\tilde{M}$ by

$$
f(q)=\exp_{\tilde{p}}\circ\,i\circ\exp_p^{-1}(q),\qquad q\in V.
$$

For all $q\in V$ there exists a unique normalized geodesic $\gamma:[0,t]\to M$
with $\gamma(0)=p$, $\gamma(t)=q$. Denote by $P_t$ the parallel transport along
$\gamma$ from $\gamma(0)$ to $\gamma(t)$. Define $\phi_t:T_q(M)\to T_{f(q)}(\tilde{M})$
by

$$
\phi_t(v)=\tilde{P}_t\circ i\circ P_t^{-1}(v),\qquad v\in T_q(M),
$$

where $\tilde{P}_t$ is the parallel transport along the normalized geodesic
$\tilde{\gamma}:[0,t]\to\tilde{M}$ given by $\tilde{\gamma}(0)=\tilde{p}$,
$\tilde{\gamma}'(0)=i(\gamma'(0))$. Denote by $R$ and $\tilde{R}$ the curvatures of
$M$ and $\tilde{M}$, respectively.

\begin{lemma}[frame-coefficient Jacobi transfer: the analytic heart of Cartan's theorem]
  \label{lem:dc-ch8-2-1-jacobi-transfer}
  \lean{Riemannian.Jacobi.jacobiFrameTransfer}
  \leanok
  \uses{def:dc-ch5-2-1}
  Let $\gamma:[0,\ell]\to M$ and $\tilde\gamma:[0,\ell]\to\tilde M$ be geodesics, read
  in fixed charts as the coordinate curves $u$ and $\tilde u$. Let $e_1(t),\dots,e_n(t)$
  be a parallel orthonormal frame along $\gamma$ and $\tilde e_1(t),\dots,\tilde e_n(t)$
  a parallel orthonormal frame along $\tilde\gamma$ ($n=\dim M$), with $e_n=u'=\gamma'$,
  $\tilde e_n=\tilde u'=\tilde\gamma'$, and let $y_1,\dots,y_n:\mathbb R\to\mathbb R$ be
  scalar coefficient functions. Suppose that $J=\sum_i y_i e_i$ satisfies the Jacobi
  equation $\tfrac{D^2 J}{dt^2}+R(u',J)u'=0$ along $\gamma$, and that at the parameter
  $t$ the two frames have matching curvature coefficients,
  $$
  \langle R(e_n,e_i)e_n,e_j\rangle=\langle\tilde R(\tilde e_n,\tilde e_i)\tilde e_n,\tilde e_j\rangle,
  \qquad i,j=1,\dots,n.
  $$
  Then the field $\tilde J=\sum_i y_i\tilde e_i$ obtained by carrying over the \emph{same}
  scalar coefficients satisfies the Jacobi equation
  $\tfrac{D^2\tilde J}{dt^2}+\tilde R(\tilde u',\tilde J)\tilde u'=0$ along $\tilde\gamma$;
  i.e. $\tilde J$ is a Jacobi field along $\tilde\gamma$.
\end{lemma}
\begin{proof}
  \leanok
  Because $e$ is parallel and orthonormal, pairing the Jacobi equation for
  $J=\sum_i y_i e_i$ with $e_j$ reduces it to do Carmo's second-order scalar system
  $y_j''+\sum_i a_{ij}\,y_i=0$, $a_{ij}=\langle R(e_n,e_i)e_n,e_j\rangle$ (the frame
  reduction: $\tfrac{D^2}{dt^2}\sum_i y_i e_i=\sum_i y_i''e_i$ since $De_i/dt=0$, and
  $R(u',\cdot)u'$ is linear). By the matching hypothesis $a_{ij}=\tilde a_{ij}$ with
  $\tilde a_{ij}=\langle\tilde R(\tilde e_n,\tilde e_i)\tilde e_n,\tilde e_j\rangle$, so
  the \emph{same} coefficients $y_i$ solve the analogous system
  $y_j''+\sum_i\tilde a_{ij}\,y_i=0$ on the $\tilde M$ side. Reassembling in the parallel
  orthonormal frame $\tilde e$ (the converse of the frame reduction: the basis expansion
  of $\tfrac{D^2\tilde J}{dt^2}+\tilde R(\tilde u',\tilde J)\tilde u'$ has these scalar
  expressions as its coefficients) gives $\tfrac{D^2\tilde J}{dt^2}+\tilde R(\tilde u',\tilde J)\tilde u'=0$.
\end{proof}

\begin{lemma}[the frame transfer in constant curvature: the matching hypothesis is automatic]
  \label{lem:dc-ch8-2-1-jacobi-transfer-const}
  \lean{Riemannian.Jacobi.jacobiFrameTransfer_isConstantCurvature}
  \leanok
  \uses{lem:dc-ch8-2-1-jacobi-transfer}
  Suppose $M$ and $\tilde M$ both have constant sectional curvature $K_0$ (the same value),
  and let $e_1,\dots,e_n$ resp.\ $\tilde e_1,\dots,\tilde e_n$ be parallel orthonormal frames
  along geodesics $\gamma$ resp.\ $\tilde\gamma$, with the velocities $u'=e_{n_0}$,
  $\tilde u'=\tilde e_{n_0}$ among the frame vectors. Then the curvature-matching hypothesis
  of \cref{lem:dc-ch8-2-1-jacobi-transfer},
  $$
  \langle R(e_i,u')u',e_j\rangle=\langle\tilde R(\tilde e_i,\tilde u')\tilde u',\tilde e_j\rangle,
  $$
  holds \emph{automatically}: by Lemma 3.4 of Chap. 4 (the pointwise constant-curvature model
  form $\langle R(x,y)z,t\rangle=K_0(\langle x,z\rangle\langle y,t\rangle-\langle y,z\rangle\langle x,t\rangle)$)
  and orthonormality, both sides equal the manifold-independent Kronecker matrix
  $K_0(\delta_{ij}-\delta_{i n_0}\delta_{n_0 j})$. Hence a Jacobi field $J=\sum_i y_i e_i$ along
  $\gamma$ carries over, with the \emph{same} scalar coefficients, to a Jacobi field
  $\tilde J=\sum_i y_i\tilde e_i$ along $\tilde\gamma$ --- with \emph{no} parallel-transport
  conjugation $\phi_t$ required. This is exactly the mechanism used by \cref{cor:dc-ch8-2-2},
  \cref{cor:dc-ch8-2-3} and \cref{thm:dc-ch8-4-1}, which only ever compare two spaces of the
  same constant curvature.
\end{lemma}
\begin{proof}
  \leanok
  Pairing the Jacobi operator $a\mapsto R(a,u')u'$ with $e_j$ and reading it in the fixed chart
  gives, via the chart $\leftrightarrow$ manifold curvature bridge and the constant-curvature
  model form of Lemma 3.4 of Chap. 4, the chart identity
  $\langle R(a,u')u',b\rangle=K_0(\langle u',u'\rangle\langle a,b\rangle-\langle a,u'\rangle\langle u',b\rangle)$
  (\lstinline{chartMetricInner\_chartCurvatureEndo\_isConstantCurvature}). Substituting
  $a=e_i$, $b=e_j$, $u'=e_{n_0}$ and using orthonormality
  ($\langle e_i,e_j\rangle=\delta_{ij}$, $\langle e_i,e_{n_0}\rangle=\delta_{i n_0}$,
  $\langle e_{n_0},e_{n_0}\rangle=1$) collapses each entry to
  $K_0(\delta_{ij}-\delta_{i n_0}\delta_{n_0 j})$, which contains no manifold data
  (\lstinline{chartCurvatureEndo\_frameCoef\_isConstantCurvature}). The two frame-coefficient
  matrices therefore coincide, discharging the hypothesis of
  \cref{lem:dc-ch8-2-1-jacobi-transfer}, whose conclusion is the claimed carry-over.
\end{proof}

\begin{lemma}[velocity-seeded parallel orthonormal frame]
  \label{lem:dc-ch8-2-1-velocity-frame}
  \lean{Riemannian.Jacobi.exists_velocitySeededParallelOrthoFrame}
  \leanok
  \uses{def:dc-ch5-2-1}
  Let $\gamma:[a,b]\to M$ be a geodesic, read in a fixed chart as the coordinate curve
  $u$, whose initial velocity $u'(a)=\gamma'(a)$ is a unit vector. Then there is a parallel
  orthonormal frame $e_1(t),\dots,e_n(t)$ along $\gamma$ (parallel: $De_i/dt=0$; orthonormal:
  $\langle e_i(t),e_j(t)\rangle=\delta_{ij}$ for all $t$) whose distinguished member is the
  velocity, $e_{n_0}(t)=\gamma'(t)$ for all $t\in[a,b]$. This is do Carmo's
  ``$e_1,\dots,e_{n-1},e_n=\gamma'$''.
\end{lemma}
\begin{proof}
  \leanok
  The unit initial velocity $\gamma'(a)$ extends to an orthonormal basis $\{e_i^0\}$ of
  $T_p(M)$ with $e_{n_0}^0=\gamma'(a)$ (Householder reflection of a generic orthonormal basis:
  a reflection preserves the inner product and can be chosen to send a basis vector to
  $\gamma'(a)$). Parallel transport of $\{e_i^0\}$ along $\gamma$ gives a parallel orthonormal
  frame $\{e_i(t)\}$ (parallel transport is a linear isometry). Both $e_{n_0}$ and the
  velocity $\gamma'$ are parallel fields along $\gamma$ --- the velocity because $\gamma$ is a
  geodesic ($D\gamma'/dt=0$) --- agreeing at $t=a$; by uniqueness of parallel transport
  $e_{n_0}(t)=\gamma'(t)$ for all $t$.
\end{proof}

\begin{lemma}[the parallel orthonormal frame along a geodesic, intrinsically]
  \label{lem:dc-ch8-2-1-parallel-ortho-frame}
  \lean{Riemannian.Jacobi.exists_parallelOrthoFrameAlongOn}
  \leanok
  \uses{def:dc-ch5-2-1}
  Let $\gamma:[a,b]\to M$ be a geodesic and let $\{e^0_i\}$ be an orthonormal family in
  $T_{\gamma(a)}M$. Then there is a family $\{e_i(t)\}$ of vector fields along $\gamma$ with
  $e_i(a)=e^0_i$, each $e_i$ parallel along $\gamma$, and
  $\langle e_i(t),e_j(t)\rangle_{\gamma(t)}=\delta_{ij}$ for \emph{every} $t\in[a,b]$ --- not
  merely at $t=a$.

  This is the manifold-level form of \cref{lem:dc-ch8-2-1-velocity-frame}'s frame: being
  chart-local, it applies to a geodesic that leaves any single chart, which is the situation
  of the antipodal argument of \cref{thm:dc-ch8-4-1}.
\end{lemma}
\begin{proof}
  \leanok
  Parallel-transport each $e^0_i$ along $\gamma$; this exists with any prescribed initial
  value. Parallel transport is an isometry --- the pairing $t\mapsto\langle e_i(t),e_j(t)
  \rangle_{\gamma(t)}$ of two parallel fields is constant along $\gamma$ --- so the Gram
  matrix at time $t$ equals the Gram matrix at time $a$, which is $\delta_{ij}$ by
  hypothesis. $\square$
\end{proof}

\begin{lemma}[the frame transported by a linear isometry: do Carmo's $\tilde e_j=\phi_t(e_j)$]
  \label{lem:dc-ch8-2-1-transported-frame}
  \lean{Riemannian.Jacobi.exists_transportedParallelOrthoFrame,
    Riemannian.Jacobi.exists_transportedParallelOrthoFrame_pair}
  \leanok
  \uses{lem:dc-ch8-2-1-parallel-ortho-frame}
  Let $i:T_pM\to T_{\tilde p}\tilde M$ be a linear isometry, let $\{e^0_j\}$ be an orthonormal
  family at $p=\gamma(a)$, and let $\tilde\gamma:[a,b]\to\tilde M$ be a geodesic with
  $\tilde\gamma(a)=\tilde p$. Then there is a family $\{\tilde e_j(t)\}$ along $\tilde\gamma$
  with $\tilde e_j(a)=i(e^0_j)$, each $\tilde e_j$ parallel along $\tilde\gamma$, and
  $\langle\tilde e_j(t),\tilde e_k(t)\rangle_{\tilde\gamma(t)}=\delta_{jk}$ for every
  $t\in[a,b]$. Taken together with \cref{lem:dc-ch8-2-1-parallel-ortho-frame} this produces
  two frames matched at time $a$ by $\tilde e_j(a)=i(e_j(a))$, which is do Carmo's
  $\tilde e_j(t)=\phi_t(e_j(t))$ at $t=a$, where $\phi_a=i$.

  Here $\phi_t=\tilde P_t\circ i\circ P_t^{-1}$, with $P_t$, $\tilde P_t$ parallel transport
  from time $a$ along $\gamma$, $\tilde\gamma$ (\cref{lem:dc-ch8-2-1-parallel-transport-map}).
  The family is the $\phi_t$-image of the frame $e$, and $\phi_t$ need never be formed in order
  to build it: since $e$ is parallel with $e_j(a)=e^0_j$ we have $P_t^{-1}(e_j(t))=e^0_j$, hence
  $\phi_t(e_j(t))=\tilde P_t(i(e^0_j))$, the parallel field on $\tilde M$ seeded by $i(e^0_j)$.
  Seeding the transport therefore delivers the frame directly. The identification
  $\tilde e_j(t)=\phi_t(e_j(t))$ itself is \cref{lem:dc-ch8-2-1-phi}.

  One gap separates this from do Carmo's frame pair, and is left to later work rather than
  assumed here: $\tilde\gamma$ is an arbitrary geodesic issuing from $\tilde p$, whereas
  do Carmo additionally requires $\tilde\gamma'(a)=i(\gamma'(a))$, and his frames carry the
  distinguished members $e_n=\gamma'$ and $\tilde e_n=\tilde\gamma'$
  (cf.\ \cref{lem:dc-ch8-2-1-velocity-frame}), neither of which is concluded here. What the
  seed buys is what the argument of \cref{thm:dc-ch8-2-1} consumes: orthonormality of
  $\tilde e$ at every time, matched to $e$ at time $a$ through $i$.
\end{lemma}
\begin{proof}
  \leanok
  Since $i$ is a linear isometry, $\langle i(e^0_j),i(e^0_k)\rangle_{\tilde p}
  =\langle e^0_j,e^0_k\rangle_p=\delta_{jk}$: the transported seed $\{i(e^0_j)\}$ is
  orthonormal at $\tilde p$. Apply \cref{lem:dc-ch8-2-1-parallel-ortho-frame} on $\tilde M$
  along $\tilde\gamma$ with that seed. $\square$
\end{proof}

\begin{lemma}[uniqueness of parallel transport, and the transport map $P_t$]
  \label{lem:dc-ch8-2-1-parallel-transport-map}
  \lean{Riemannian.Jacobi.IsParallelFieldAlongOn.eqOn_of_initial,
    Riemannian.Jacobi.parallelTransportAlong,
    Riemannian.Jacobi.metricInner_parallelTransportAlong,
    Riemannian.Jacobi.parallelTransportAlongEquiv}
  \leanok
  \uses{def:dc-ch5-2-1}
  Let $\gamma:[a,b]\to M$ be a geodesic. Two fields parallel along $\gamma$ that agree at $a$
  agree on all of $[a,b]$. Consequently the assignment carrying $w_0\in T_{\gamma(a)}M$ to the
  value at $t$ of the parallel field along $\gamma$ seeded by $w_0$ is a well-defined map
  $$
  P_t:T_{\gamma(a)}M\longrightarrow T_{\gamma(t)}M,\qquad t\in[a,b],
  $$
  which is linear, satisfies $P_a=\mathrm{id}$, and is an isometry:
  $\langle P_tu,P_tw\rangle_{\gamma(t)}=\langle u,w\rangle_{\gamma(a)}$. In particular $P_t$ is
  injective, hence a linear isomorphism onto $T_{\gamma(t)}M$.
\end{lemma}
\begin{proof}
  \leanok
  \uses{def:dc-ch5-2-1}
  Let $v,w$ be parallel along $\gamma$ with $v(a)=w(a)$, and set $d=v-w$. The parallel
  transport equation $\frac{Dd}{dt}=0$ is linear in $d$, so $d$ is parallel along $\gamma$,
  and $d(a)=0$. The pairing of two parallel fields along $\gamma$ is constant, so
  $$
  \langle d(t),d(t)\rangle_{\gamma(t)}=\langle d(a),d(a)\rangle_{\gamma(a)}=0
  \qquad(t\in[a,b]),
  $$
  and since $g_{\gamma(t)}$ is positive definite, $d(t)=0$; that is, $v(t)=w(t)$.

  Uniqueness makes $P_t$ well defined, and existence of a parallel field with prescribed value
  at $a$ makes it total. For linearity, let $v,w$ be the parallel fields seeded by $w_0,w_1$.
  Then $v+w$ is parallel and $(v+w)(a)=w_0+w_1$, so by uniqueness $v+w$ is the parallel field
  seeded by $w_0+w_1$; evaluating at $t$ gives $P_t(w_0+w_1)=P_tw_0+P_tw_1$, and homogeneity is
  the same argument applied to $rv$. $P_a=\mathrm{id}$ is the seeding condition. The isometry
  statement is once more the constancy of the pairing of two parallel fields, applied to the
  fields seeded by $u$ and $w$. Finally, if $P_tw_0=0$ then
  $\langle w_0,w_0\rangle_{\gamma(a)}=\langle P_tw_0,P_tw_0\rangle_{\gamma(t)}=0$, so $w_0=0$;
  an injective endomorphism of a finite-dimensional space is an isomorphism. $\square$
\end{proof}

\begin{lemma}[do Carmo's $\phi_t$, and $\tilde e(t)=\phi_t(e(t))$]
  \label{lem:dc-ch8-2-1-phi}
  \lean{Riemannian.Jacobi.parallelTransportConjugate,
    Riemannian.Jacobi.parallelTransportConjugate_left,
    Riemannian.Jacobi.metricInner_parallelTransportConjugate,
    Riemannian.Jacobi.eq_parallelTransportConjugate_of_isParallelFieldAlongOn}
  \leanok
  \uses{lem:dc-ch8-2-1-parallel-transport-map}
  Let $\gamma:[a,b]\to M$ and $\tilde\gamma:[a,b]\to\tilde M$ be geodesics and let
  $i:T_{\gamma(a)}M\to T_{\tilde\gamma(a)}\tilde M$ be linear. For $t\in[a,b]$ put
  $$
  \phi_t=\tilde P_t\circ i\circ P_t^{-1}:
  T_{\gamma(t)}M\longrightarrow T_{\tilde\gamma(t)}\tilde M,
  $$
  with $P_t$, $\tilde P_t$ the parallel transports of
  \cref{lem:dc-ch8-2-1-parallel-transport-map}. Then $\phi_a=i$; if $i$ is a linear isometry
  then so is every $\phi_t$; and if $v$ is parallel along $\gamma$ and $\tilde v$ is parallel
  along $\tilde\gamma$ with $\tilde v(a)=i(v(a))$, then
  $$
  \tilde v(t)=\phi_t(v(t))\qquad\text{for every }t\in[a,b].
  $$
\end{lemma}
\begin{proof}
  \leanok
  \uses{lem:dc-ch8-2-1-parallel-transport-map}
  $\phi_a=i$ because $P_a$ and $\tilde P_a$ are the identity. If $i$ is a linear isometry then
  $\phi_t$ is a composite of three isometries, $P_t^{-1}$ and $\tilde P_t$ being isometries by
  \cref{lem:dc-ch8-2-1-parallel-transport-map}; hence $\phi_t$ is one.

  For the last claim, $v$ is parallel with value $v(a)$ at $a$, so by the uniqueness clause of
  \cref{lem:dc-ch8-2-1-parallel-transport-map} it is the parallel field seeded by $v(a)$, i.e.
  $P_t(v(a))=v(t)$, whence $P_t^{-1}(v(t))=v(a)$. Therefore
  $$
  \phi_t(v(t))=\tilde P_t\big(i(v(a))\big)=\tilde P_t\big(\tilde v(a)\big)=\tilde v(t),
  $$
  the last equality because $\tilde v$ is parallel along $\tilde\gamma$ with value $\tilde v(a)$
  at $a$, hence is the parallel field seeded by $\tilde v(a)$. $\square$
\end{proof}

\begin{lemma}[the chart frame coefficient is the intrinsic curvature form]
  \label{lem:dc-ch8-2-1-curvature-bridge}
  \lean{Riemannian.Jacobi.chartMetricInner_chartCurvatureEndo_eq_curvatureFormAt,
    Riemannian.Jacobi.chartFrameRealize_tangentCoordChange,
    Riemannian.Jacobi.chartMetricInner_chartCurvatureEndo_chartVectorRep_eq_curvatureFormAt}
  \leanok
  For $v,a,b\in T_qM$, the frame coefficient $\langle\mathcal R(a,v)v,b\rangle$ computed in a
  chart --- the form in which the Jacobi equation is actually run --- equals the intrinsic
  curvature form $R(v,a,v,b)$ evaluated on $v,a,v,b$ themselves. No curvature hypothesis is
  assumed.
\end{lemma}
\begin{proof}
  \leanok
  Pairing the chart Jacobi operator against a chart vector gives
  $\langle\mathcal R(a,v)v,b\rangle$. The manifold--chart curvature bridge identifies
  $R(\hat a,\hat v,\hat v,\hat b)$ --- the intrinsic form on the chart-frame realizations
  $\hat\cdot$ of the chart vectors --- with the \emph{negative} of that quantity, and
  antisymmetry of $R$ in its first pair rewrites $R(\hat a,\hat v,\hat v,\hat b)$ as
  $-R(\hat v,\hat a,\hat v,\hat b)$; the two signs cancel. Finally, realizing the chart
  reading of a tangent vector in the chart frame at its own foot returns the vector
  (the readback inverse to the chart reading), so for intrinsic $v,a,b$ no realizations
  remain. $\square$
\end{proof}

\begin{lemma}[E.\ Cartan's curvature hypothesis, converted to the frame coefficients]
  \label{lem:dc-ch8-2-1-hmatch-transfer}
  \lean{Riemannian.Jacobi.chartMetricInner_chartCurvatureEndo_transfer_of_curvatureFormAt}
  \leanok
  \uses{lem:dc-ch8-2-1-curvature-bridge}
  Let $\varphi:T_qM\to T_{\tilde q}\tilde M$ be any map under which the curvature forms
  correspond,
  $$
  \langle R(x,y)z,w\rangle=\big\langle\tilde R(\varphi x,\varphi y)\varphi z,\varphi w\big\rangle
  \qquad\text{for all }x,y,z,w\in T_qM ,
  $$
  which is the hypothesis of \cref{thm:dc-ch8-2-1} with $\varphi=\phi_t$. Then the chart frame
  coefficients of $M$ and $\tilde M$ agree: for all $v,a,b\in T_qM$,
  $$
  \langle\mathcal R(a,v)v,b\rangle=\big\langle\tilde{\mathcal R}(\varphi a,\varphi v)\varphi v,
  \varphi b\big\rangle .
  $$
  Taking $v=\gamma'$ and $a,b$ frame vectors, this is the intrinsic content of the matching
  hypothesis of \cref{lem:dc-ch8-2-1-jacobi-transfer} and of
  \cref{lem:dc-ch8-2-1-jacobi-norm-transfer-general} --- the single point at which the
  curvature hypothesis of \cref{thm:dc-ch8-2-1} is consumed.

  It is not yet that matching hypothesis itself. Those lemmas read
  their frames in a fixed chart along a coordinate curve, so identifying the two demands three
  further things, none of them about curvature: that $\varphi(\gamma'(t))=\tilde\gamma'(t)$
  (which needs $\tilde\gamma'(a)=i(\gamma'(a))$); that its frame vectors are the chart
  readings of the present intrinsic ones under the present $\varphi$; and that its
  parallelism hypothesis, stated in one fixed chart, follows from parallelism along
  $\gamma$ in the sense of \cref{lem:dc-ch8-2-1-parallel-ortho-frame}, which is chart-local.
  That interface is the unbuilt residual of \cref{thm:dc-ch8-2-1}.

  Note $\varphi$ is arbitrary: no linearity, continuity or isometry is used, and no
  parallel-transport theory enters --- $\varphi$ is data, supplied by \cref{thm:dc-ch8-2-1}'s
  own hypothesis as $\phi_t$. The additional structure of $\phi_t$ is what makes the
  transported frame orthonormal (\cref{lem:dc-ch8-2-1-transported-frame}), which is a separate
  requirement.
\end{lemma}
\begin{proof}
  \leanok
  Rewrite each side by \cref{lem:dc-ch8-2-1-curvature-bridge} into the intrinsic curvature
  forms $R(v,a,v,b)$ and $\tilde R(\varphi v,\varphi a,\varphi v,\varphi b)$, and apply the
  hypothesis at $(v,a,v,b)$. $\square$
\end{proof}

\begin{lemma}[the Jacobi field read in a parallel orthonormal frame]
  \label{lem:dc-ch8-2-1-frame-coeff-bridge}
  \lean{Riemannian.Jacobi.isJacobiPairOn_of_isJacobiFieldOn}
  \leanok
  \uses{def:dc-ch5-2-1,lem:dc-ch8-2-1-parallel-ortho-frame}
  Let $\gamma$ be a geodesic of $M$, read in a fixed chart as the coordinate curve $u$
  staying in the chart domain, and let the pair $(J,DJ)$ satisfy the Jacobi pair system
  along $u$ on a parameter window $[a',b']$,
  $$
  \tfrac{D}{dt}J=DJ,\qquad \tfrac{D}{dt}DJ=-\mathcal R(J,u')u' ,
  $$
  where $\langle\cdot,\cdot\rangle$ and $\mathcal R$ denote the metric and curvature read in
  that chart. Let $e_1(t),\dots,e_n(t)$ ($n=\dim M$) be a parallel orthonormal frame along
  $u$ in the same chart. For a parameter interval $[a,b]\subset(a',b')$ set
  $$
  F_i(t)=\langle J(t),e_i(t)\rangle,\qquad V_i(t)=\langle DJ(t),e_i(t)\rangle .
  $$
  Then on $[a,b]$ the pair $(F,V)$ solves the first-order system
  $$
  F'=V,\qquad V'=-A(t)F,\qquad A(t)_{ij}=a_{ij}(t)=\langle\mathcal R(e_i,u')u',e_j\rangle ,
  $$
  i.e.\ $F_j''+\sum_i a_{ij}F_i=0$: do Carmo's scalar system for $J=\sum_i F_ie_i$.

  Taking $[a,b]$ strictly inside the window $[a',b']$ makes every parameter of $[a,b]$ an
  interior one, where the derivatives of the pair system are two-sided; this is what the
  product rule below requires.
\end{lemma}
\begin{proof}
  \leanok
  Both halves are metric compatibility read against a parallel frame. Since $De_j/dt=0$,
  the product rule $\tfrac{d}{dt}\langle X,e_j\rangle=\langle DX/dt,e_j\rangle+\langle
  X,De_j/dt\rangle$ reduces to
  $$
  \tfrac{d}{dt}\langle X,e_j\rangle=\langle DX/dt,e_j\rangle
  $$
  for any field $X$ along $\gamma$. Applying this to $X=J$ and using $DJ/dt=DJ$ gives
  $F_j'=\langle DJ,e_j\rangle=V_j$. Applying it to $X=DJ$ and using the Jacobi equation
  $D(DJ)/dt=-R(\gamma',J)\gamma'$ gives
  $V_j'=-\langle R(\gamma',J)\gamma',e_j\rangle$. Expanding $J=\sum_i F_ie_i$
  (orthonormality and completeness of the frame) and using linearity of
  $x\mapsto R(\gamma',x)\gamma'$,
  $$
  \langle R(\gamma',J)\gamma',e_j\rangle=\sum_i F_i\langle R(\gamma',e_i)\gamma',e_j\rangle
  =\sum_i a_{ij}F_i=(A(t)F)_j ,
  $$
  so $V_j'=-(A(t)F)_j$. $\square$
\end{proof}

\begin{lemma}[transfer of the Jacobi norm under matching curvature coefficients]
  \label{lem:dc-ch8-2-1-jacobi-norm-transfer-general}
  \lean{Riemannian.Jacobi.chartMetricInner_transfer_of_jacobiCoef_match}
  \leanok
  \uses{lem:dc-ch8-2-1-frame-coeff-bridge,lem:dc-ch8-2-1-hmatch-transfer,def:dc-ch5-2-1}
  Let $\gamma$ in $M$ and $\tilde\gamma$ in $\tilde M$ ($\dim M=\dim\tilde M=n$) be geodesics
  read in fixed charts as the coordinate curves $u$ and $\tilde u$, each staying in its chart
  domain, let $(J,DJ)$ and $(\tilde J,D\tilde J)$ satisfy the Jacobi pair systems along $u$
  and $\tilde u$ on a window $[a',b']$, and let $e$ and $\tilde e$ be parallel orthonormal
  frames along $u$ and $\tilde u$ in those charts; as in
  \cref{lem:dc-ch8-2-1-frame-coeff-bridge}, $\langle\cdot,\cdot\rangle$ and $\mathcal R$
  denote the metric and curvature read in those charts. Let $[a,b]\subset(a',b')$, and
  suppose that on $[a,b]$ the chart curvature operator $w\mapsto\mathcal R(w,u')u'$ along
  $u$, the frame vectors $e_i$, and the chart metric coefficients along $u$ are all
  continuous, and that the curvature coefficients agree there,
  $$
  \langle\mathcal R(e_i,u')u',e_j\rangle
  =\big\langle\tilde{\mathcal R}(\tilde e_i,\tilde u')\tilde u',\tilde e_j\big\rangle
  \qquad\text{for all }i,j\text{ and all }t\in[a,b],
  $$
  and that the two fields have the same frame data at $a$,
  $\langle J(a),e_i(a)\rangle=\langle\tilde J(a),\tilde e_i(a)\rangle$ and
  $\langle DJ(a),e_i(a)\rangle=\langle D\tilde J(a),\tilde e_i(a)\rangle$ for all $i$.
  Then
  $$
  |\tilde J(t)|^2=|J(t)|^2\qquad\text{for all }t\in[a,b] .
  $$
\end{lemma}
\begin{proof}
  \leanok
  By \cref{lem:dc-ch8-2-1-frame-coeff-bridge} the frame coefficients $(F,V)$ of $(J,DJ)$
  solve $F'=V$, $V'=-A(t)F$ on $[a,b]$ with
  $A(t)_{ij}=\langle\mathcal R(e_i,u')u',e_j\rangle$, and likewise $(\tilde F,\tilde V)$
  solve $\tilde F'=\tilde V$, $\tilde V'=-\tilde A(t)\tilde F$ with
  $\tilde A(t)_{ij}=\langle\tilde{\mathcal R}(\tilde e_i,\tilde u')\tilde u',\tilde e_j\rangle$.
  The matching hypothesis says precisely $A(t)=\tilde A(t)$ for $t\in[a,b]$; since the system
  only reads its coefficient there, $(F,V)$ and $(\tilde F,\tilde V)$ solve the \emph{same}
  linear system on $[a,b]$. The entries $a_{ij}=\langle\mathcal R(e_i,u')u',e_j\rangle$ are
  built from the chart curvature operator along $u$, the frame vectors and the chart metric
  coefficients along $u$, so by the continuity hypotheses $A$ is continuous on $[a,b]$, and
  uniqueness for this system (\cref{def:dc-ch5-2-1}) applies: two solutions sharing data at
  $a$ agree, so from $F(a)=\tilde F(a)$ and $V(a)=\tilde V(a)$ we get $F=\tilde F$ on
  $[a,b]$.

  Finally, orthonormality of each frame reads the norm off the coefficients: expanding
  $J=\sum_i F_ie_i$ and $\tilde J=\sum_i\tilde F_i\tilde e_i$,
  $$
  |J(t)|^2=\sum_i F_i(t)^2,\qquad |\tilde J(t)|^2=\sum_i\tilde F_i(t)^2 ,
  $$
  and these agree because $F=\tilde F$ on $[a,b]$. $\square$

  Note that no closed-form solution of the scalar system is used, and no hypothesis on the
  sign or constancy of the curvature: only that the two coefficient matrices coincide. This
  is the general form of the comparison that
  \cref{lem:dc-ch8-2-1-jacobi-norm-transfer} performs, for constant curvature, by solving
  the system explicitly.
\end{proof}

\begin{lemma}[the transported frame and velocity, read in a fixed chart]
  \label{lem:dc-ch8-2-1-frame-chart-interface}
  \lean{Riemannian.Jacobi.jacobiCoef_match_of_curvatureFormAt,
    Riemannian.Jacobi.jacobiCoef_match_of_curvatureFormAt_parallelTransportConjugate}
  \leanok
  \uses{lem:dc-ch8-2-1-hmatch-transfer,lem:dc-ch8-2-1-transported-frame,
    lem:dc-ch8-2-1-parallel-transport-map,lem:dc-ch8-2-1-phi}
  Let $\gamma:[a,b]\to M$ and $\tilde\gamma:[a,b]\to\tilde M$ be geodesics, each of which
  \emph{stays in the source of one fixed chart}, read there as the coordinate curves $u$ and
  $\tilde u$. Let $e_1,\dots,e_n$ be parallel fields along $\gamma$ and $\tilde e_1,\dots,
  \tilde e_n$ parallel fields along $\tilde\gamma$ with matched seeds
  $\tilde e_j(a)=i(e_j(a))$, and suppose $\tilde\gamma'(a)=i(\gamma'(a))$; read all of them
  in those same charts. Then the curvature hypothesis of \cref{thm:dc-ch8-2-1} implies the
  matching hypothesis of \cref{lem:dc-ch8-2-1-jacobi-norm-transfer-general} in the form that
  lemma requires it, namely between the \emph{chart} frame coefficients:
  $$
  \langle\mathcal R(e_i,u')u',e_j\rangle
  =\big\langle\tilde{\mathcal R}(\tilde e_i,\tilde u')\tilde u',\tilde e_j\big\rangle
  \qquad\text{for all }i,j\text{ and all }t\in[a,b] .
  $$
  More generally the same conclusion holds with $\phi_t$ replaced by \emph{any} family of
  maps $T_{\gamma(t)}M\to T_{\tilde\gamma(t)}\tilde M$ under which the curvature forms
  correspond and which carries $\gamma'(t)$ to $\tilde\gamma'(t)$ and $e_j(t)$ to
  $\tilde e_j(t)$: no linearity, continuity or isometry of the family is used, and no
  parallel-transport theory enters. The two Lean statements are that general form and its
  instantiation at $\phi_t$.

  Two hypotheses that might be expected are \emph{absent}. The frames need not be
  orthonormal, and need not contain the velocity: do Carmo writes the coefficients as
  $\langle R(e_n,e_i)e_n,e_j\rangle$ with $e_n=\gamma'$, but the coefficient of
  \cref{lem:dc-ch8-2-1-jacobi-norm-transfer-general} contracts the curvature against $u'$
  directly and never against a frame member, so $e_n=\gamma'$ is not needed here.
  (Orthonormality is still needed by that lemma itself, to read the norm off the
  coefficients.)
\end{lemma}
\begin{proof}
  \leanok
  First, $\phi_t$ carries the velocity to the velocity:
  $\phi_t(\gamma'(t))=\tilde\gamma'(t)$. Indeed a geodesic has parallel velocity, so
  $\gamma'(t)=P_t(\gamma'(a))$ and $\tilde\gamma'(t)=\tilde P_t(\tilde\gamma'(a))$, where
  $P_t,\tilde P_t$ denote parallel transport from $a$ to $t$
  (\cref{lem:dc-ch8-2-1-parallel-transport-map}). Since
  $\phi_t=\tilde P_t\circ i\circ P_t^{-1}$ (\cref{lem:dc-ch8-2-1-phi}),
  $$
  \phi_t(\gamma'(t))=\tilde P_t\big(i(P_t^{-1}\gamma'(t))\big)=\tilde P_t\big(i(\gamma'(a))\big)
  =\tilde P_t(\tilde\gamma'(a))=\tilde\gamma'(t) ,
  $$
  using $\tilde\gamma'(a)=i(\gamma'(a))$. (Both $\gamma'$ and $\tilde\gamma'$ being parallel
  with $i$-matched seeds, this is \cref{lem:dc-ch8-2-1-transported-frame} applied to the
  velocity fields; $P_t$ itself never needs forming.)

  Second, $\tilde e_j(t)=\phi_t(e_j(t))$ by \cref{lem:dc-ch8-2-1-transported-frame}.

  Hence applying \cref{lem:dc-ch8-2-1-hmatch-transfer} with $\varphi=\phi_t$ at
  $(v,a,b)=(\gamma'(t),e_i(t),e_j(t))$ gives the matching of the chart frame coefficients of
  the \emph{chart readings} of $\gamma'(t),e_i(t),e_j(t)$ and of their images under $\phi_t$,
  which by the two paragraphs above are the chart readings of
  $\tilde\gamma'(t),\tilde e_i(t),\tilde e_j(t)$. $\square$

  The three identifications this proof glosses over — which were, until they were carried
  out, the residual of \cref{thm:dc-ch8-2-1} — turned out to be cheap, and it is worth
  recording why, since each was resolved by a choice rather than by an argument. First,
  \cref{lem:dc-ch8-2-1-hmatch-transfer} presents its vectors as the chart readings, at their
  own foot, of intrinsic tangent vectors, whereas
  \cref{lem:dc-ch8-2-1-jacobi-norm-transfer-general} presents $e_i$, $u'$ as coordinate
  fields along the coordinate curve in one fixed chart: taking the latter to \emph{be} the
  own-foot chart readings of the former makes the identification hold definitionally.
  Second, the curvature operator of that lemma is read along the coordinate curve and had to
  be matched with the pointwise chart curvature of \cref{lem:dc-ch8-2-1-hmatch-transfer};
  the two are both contractions of the same chart curvature tensor, so this is one
  computation, valid wherever the curve reads into the interior of the chart target. Third,
  the chart velocity $u'$ is, at a point of the chart source, exactly the own-foot chart
  reading of the intrinsic velocity $\gamma'(t)$ — a fact already available from Ch.~5.
  Parallelism needed no interfacing at all: the statement quantifies over parallel
  \emph{fields}, and their chart readings enter only through $\phi_t$.
\end{proof}

\begin{lemma}[the Cartan Jacobi norm transfer in variable curvature, intrinsically]
  \label{lem:dc-ch8-2-1-jacobi-norm-transfer-manifold}
  \lean{Riemannian.Jacobi.metricInner_jacobiField_transfer_of_curvatureFormAt}
  \leanok
  \uses{lem:dc-ch8-2-1-frame-chart-interface,
    lem:dc-ch8-2-1-jacobi-norm-transfer-general,
    lem:dc-ch8-2-1-parallel-ortho-frame,def:dc-ch5-2-1}
  Let $\gamma,\tilde\gamma:[a',b']\to M,\tilde M$ be geodesics, each staying in the source of
  one fixed chart; let $e_1,\dots,e_n$ and $\tilde e_1,\dots,\tilde e_n$ be parallel
  \emph{orthonormal} frames along them; and let $\phi_t$ carry $\gamma'(t)$ to
  $\tilde\gamma'(t)$, $e_j(t)$ to $\tilde e_j(t)$, and the curvature form of $M$ at
  $\gamma(t)$ to that of $\tilde M$ at $\tilde\gamma(t)$ (E. Cartan's hypothesis). Let $J$ be
  a Jacobi field along $\gamma$ and $\tilde J$ one along $\tilde\gamma$, and let
  $[a,b]\subset(a',b')$. If the two have the same frame data at $a$, then
  $$
  |\tilde J(t)|=|J(t)|\qquad\text{for all }t\in[a,b] .
  $$
  This is \cref{lem:dc-ch8-2-1-jacobi-norm-transfer-general} stated between \emph{intrinsic}
  Jacobi fields rather than their chart readings, and it is the variable-curvature
  counterpart of \cref{lem:dc-ch8-2-1-jacobi-norm-transfer}: no closed-form solution, and no
  hypothesis on the sign or constancy of the curvature.
\end{lemma}
\begin{proof}
  \leanok
  Localize both Jacobi fields, both frames and both geodesics into their fixed charts, and
  apply \cref{lem:dc-ch8-2-1-jacobi-norm-transfer-general}. The chart hypotheses that lemma
  asks for are all consequences of the manifold data: a manifold Jacobi field restricts to a
  chart certificate on any window whose $\gamma$-image lies in one chart source, and likewise
  for a parallel field, whose chart certificate gives $\nabla e_j=0$ in that chart at every
  \emph{interior} time; the chart reading is norm-faithful, which turns intrinsic
  orthonormality into the chart-Gram orthonormality required; and the chart velocity and the
  chart curvature operator of a geodesic are continuous. Its matching hypothesis is
  \cref{lem:dc-ch8-2-1-frame-chart-interface}. Finally the conclusion, a chart-Gram norm, is
  read back as the intrinsic norm by the same norm-faithfulness. $\square$

  The restriction $[a,b]\subset(a',b')$ is not cosmetic: $\nabla e_j=0$ in a fixed chart is a
  two-sided derivative statement whereas a parallel field's chart certificate is one-sided,
  so the frames must be run on an interval strictly larger than the one the conclusion is
  read on.
\end{proof}

\begin{lemma}[time reversal of the Jacobi system]
  \label{lem:dc-ch8-2-1-jacobi-reversal}
  \lean{Riemannian.Jacobi.IsJacobiFieldAlongOn.comp_neg}
  \leanok
  \uses{def:dc-ch5-2-1}
  Let $\gamma$ be a geodesic on $[a,b]$ and let $(J,DJ)$ be a Jacobi pair along it. Then
  $t\mapsto\gamma(-t)$ is a geodesic on $[-b,-a]$, and
  $$
  J^-(t)=J(-t),\qquad DJ^-(t)=-DJ(-t)
  $$
  is a Jacobi pair along it. The sign on $DJ$ is not a convention: the covariant derivative
  is a velocity, and reversing time reverses it.
\end{lemma}
\begin{proof}
  \leanok
  Read in a chart, the Jacobi pair system is
  $$
  J'=DJ-\Gamma(u',J),\qquad DJ'=-R(u;J,u',u')-\Gamma(u',DJ),
  $$
  where $u$ is the chart reading of $\gamma$. Reversing time flips the chart velocity,
  $(u\circ{-})'(t)=-u'(-t)$, and leaves $u$ itself alone. Since $\Gamma$ is \emph{linear in
  its first slot} — this is the contraction $\Gamma(v,w)$ with the field in the second slot,
  not the quadratic $\Gamma(v,v)$ of the geodesic equation — we get
  $$
  (J^-)'(t)=-J'(-t)=-DJ(-t)+\Gamma(u'(-t),J(-t))=DJ^-(t)-\Gamma((u\circ{-})'(t),J^-(t)),
  $$
  which is the first equation for the reversed pair. For the second, the curvature term is
  \emph{unchanged}, because its two velocity slots each flip sign and the two flips cancel;
  and $\Gamma((u\circ{-})',DJ^-)=\Gamma(-u'(-t),-DJ(-t))=\Gamma(u'(-t),DJ(-t))$ by linearity
  in both slots. Hence
  $$
  (DJ^-)'(t)=DJ'(-t)=-R(u;J,u',u')(-t)-\Gamma(u',DJ)(-t)
  $$
  is exactly the second equation for the reversed pair. $\square$
\end{proof}

\begin{lemma}[backward uniqueness of Jacobi fields]
  \label{lem:dc-ch8-2-1-jacobi-backward-uniqueness}
  \lean{Riemannian.Jacobi.IsJacobiFieldAlongOn.eqOn_zero_of_right}
  \leanok
  \uses{lem:dc-ch8-2-1-jacobi-reversal,def:dc-ch5-2-1}
  A Jacobi field along a geodesic on $[a,b]$ with $J(b)=0$ and $DJ(b)=0$ vanishes identically
  on $[a,b]$.
\end{lemma}
\begin{proof}
  \leanok
  Apply \cref{lem:dc-ch8-2-1-jacobi-reversal}: the reversed pair is a Jacobi field along the
  reversed geodesic on $[-b,-a]$, whose \emph{left} endpoint $-b$ is the original's right
  endpoint, and its data there is $(J(b),-DJ(b))=(0,0)$. The forward Grönwall uniqueness for
  Jacobi fields then makes it vanish on $[-b,-a]$; read back at $-t$. $\square$
\end{proof}

\begin{lemma}[a Jacobi field with prescribed data at an interior time]
  \label{lem:dc-ch8-2-1-jacobi-interior-data}
  \lean{Riemannian.Jacobi.exists_isJacobiFieldAlongOn_at}
  \leanok
  \uses{lem:dc-ch8-2-1-jacobi-backward-uniqueness,def:dc-ch5-2-1}
  Let $\gamma$ be a geodesic on $[a,b]$, let $t_0\in[a,b]$, and let $J_0,DJ_0\in
  T_{\gamma(t_0)}M$. Then there is a Jacobi field $(J,DJ)$ along $\gamma$ \emph{on all of
  $[a,b]$} with $J(t_0)=J_0$ and $DJ(t_0)=DJ_0$.

  Existence of Jacobi fields is otherwise only available seeded at the \emph{left endpoint}
  $a$; this is the same statement seeded anywhere.
\end{lemma}
\begin{proof}
  \leanok
  Consider the propagator $\Psi:(F,V)\mapsto (J_{F,V}(t_0),DJ_{F,V}(t_0))$, where $J_{F,V}$
  is the Jacobi field on $[a,b]$ with data $(F,V)$ at the left endpoint $a$. It is linear on
  $T_{\gamma(a)}M\times T_{\gamma(a)}M$, by superposition together with uniqueness of the
  field with given left-endpoint data.

  $\Psi$ is \emph{injective}: if $\Psi(F,V)=0$ then $J_{F,V}$ restricted to $[a,t_0]$ is a
  Jacobi field vanishing with its covariant derivative at the \emph{right} endpoint $t_0$, so
  $J_{F,V}$ vanishes on $[a,t_0]$ by
  \cref{lem:dc-ch8-2-1-jacobi-backward-uniqueness}, and in particular $(F,V)=0$ (for $t_0=a$
  there is nothing to prove). Injectivity is exactly where backward uniqueness — and hence
  the time reversal — is needed; forward uniqueness alone would not give it.

  Since $\gamma$ is a curve in an $n$-manifold, the domain and codomain of $\Psi$ both have
  dimension $2n$, so an injective $\Psi$ is surjective. Pull $(J_0,DJ_0)$ back through
  $\Psi$ and take the corresponding field. $\square$
\end{proof}

\begin{lemma}[from the norm transfer to the local isometry, in variable curvature]
  \label{lem:dc-ch8-2-1-exp-norm-transfer-general}
  \lean{Riemannian.Jacobi.metricInner_mfderiv_eq_of_semiconjugacy_of_curvatureFormAt}
  \leanok
  \uses{lem:dc-ch8-2-1-jacobi-norm-transfer-manifold,cor:dc-ch5-2-5,
    lem:dc-ch8-2-1-jacobi-interior-data,
    lem:dc-ch8-2-1-exp-norm-transfer,lem:dc-ch8-2-1-polarization,
    lem:dc-ch8-2-1-exp-equiv,lem:dc-ch8-2-1-mfderiv-bridge,
    lem:dc-ch8-2-2-semiconjugacy,lem:dc-ch8-2-1-no-conjugate}
  Let $M,\tilde M$ be complete, $p\in M$, $\tilde p\in\tilde M$, let $i:T_pM\to
  T_{\tilde p}\tilde M$ be a linear isometry, and let $v\in T_pM$. Write $\gamma=\gamma_v$
  and $\tilde\gamma=\gamma_{i(v)}$ for the two geodesics, and $q=\exp_p(v)$. Suppose:
  \begin{itemize}
    \item \emph{(single chart)} for some window $[a',b']$ with $a'<0<1<b'$, each of $\gamma$,
      $\tilde\gamma$ stays in the source of \emph{one} chart on $[a',b']$;
    \item \emph{(frames)} $e_1,\dots,e_n$ and $\tilde e_1,\dots,\tilde e_n$ are parallel
      \emph{orthonormal} frames along $\gamma$, $\tilde\gamma$ on $[a',b']$;
    \item \emph{(E. Cartan's hypothesis)} $\phi_t$ carries $\gamma'(t)$ to
      $\tilde\gamma'(t)$, $e_j(t)$ to $\tilde e_j(t)$, and the curvature form of $M$ at
      $\gamma(t)$ to that of $\tilde M$ at $\tilde\gamma(t)$, for $t\in[a',b']$; and
      $\phi_0=i$;
    \item \emph{(no conjugacy)} $1$ is not a conjugate parameter along $\gamma$ or along
      $\tilde\gamma$.
  \end{itemize}
  Let $f$ be \emph{any} map differentiable at $q$ that satisfies the semiconjugacy
  $f\circ\exp_p=\exp_{\tilde p}\circ\,i$ near $v$. Then $f$ preserves the metric at $q$:
  $$
  \langle u,u'\rangle_q=\langle df_q(u),df_q(u')\rangle_{f(q)}\qquad\text{for all }u,u'\in T_qM .
  $$
  This is the variable-curvature counterpart of \cref{lem:dc-ch8-2-2-semiconjugacy}, which
  proves the same conclusion from a shared constant $K_0$ rather than from E. Cartan's
  hypothesis. The single-chart and no-conjugacy clauses are \emph{hypotheses} here; see the
  discussion after the proof.
\end{lemma}
\begin{proof}
  \leanok
  The chain is the one \cref{cor:dc-ch8-2-2} runs for constant curvature
  (\cref{lem:dc-ch8-2-1-exp-norm-transfer} and its consumers), re-run against
  \cref{lem:dc-ch8-2-1-jacobi-norm-transfer-manifold}.

  Since $1$ is not conjugate along $\gamma$, $d(\exp_p)_v$ is onto
  (\cref{lem:dc-ch8-2-1-exp-equiv}), so it suffices to prove the claim on vectors
  $d(\exp_p)_v(Z)$. Realize each side as the endpoint of a Jacobi field vanishing at $0$:
  by \cref{cor:dc-ch5-2-5}, $d(\exp_p)_v(Z)=J(1)$ where $J(0)=0$, $DJ(0)=Z$, and likewise
  $d(\exp_{\tilde p})_{i(v)}(i(Z))=\tilde J(1)$ with $\tilde J(0)=0$, $D\tilde J(0)=i(Z)$.
  \emph{These fields are produced by \cref{lem:dc-ch8-2-1-jacobi-interior-data}, on the whole
  window $[a',b']$}: the norm transfer needs them there, while their data is pinned at $0$,
  which is interior — see the remark below.

  The two fields have matching frame data at $0$. For the value this is trivial, both being
  $0$. For the covariant derivative, $\tilde e_k(0)=\phi_0(e_k(0))=i(e_k(0))$, so
  $$
  \langle i(Z),\tilde e_k(0)\rangle_{\tilde p}=\langle i(Z),i(e_k(0))\rangle_{\tilde p}
    =\langle Z,e_k(0)\rangle_p
  $$
  because $i$ is a linear isometry. Note this uses only $\phi_0=i$: neither orthonormality
  of the frames nor any particular seeding of them enters here.

  So \cref{lem:dc-ch8-2-1-jacobi-norm-transfer-manifold} applies and gives
  $|\tilde J(1)|=|J(1)|$, i.e. $|d(\exp_{\tilde p})_{i(v)}(i(Z))|=|d(\exp_p)_v(Z)|$. The
  chain rule along the semiconjugacy identifies $df_q(d(\exp_p)_v(Z))$ with
  $d(\exp_{\tilde p})_{i(v)}(i(Z))$ (\cref{lem:dc-ch8-2-1-mfderiv-bridge}), which turns that
  into $|df_q(w)|=|w|$ for every $w\in T_qM$; polarization
  (\cref{lem:dc-ch8-2-1-polarization}) upgrades the diagonal to the full inner product.
  $\square$

  \textbf{Why the interior seed.} The constant-curvature route reads its two Jacobi fields on
  $[0,1]$ and seeds them at the \emph{left endpoint} $0$, which is what the existence theorem
  for Jacobi fields directly supplies. The variable-curvature transfer cannot: its parallel
  frames carry a two-sided chart-flatness certificate, so it needs the fields on a window
  $[a',b']$ with $[0,1]\subset(a',b')$, while \cref{cor:dc-ch5-2-5} still pins them by their
  data at $0$ — now an interior time. That is the one genuinely new analytic input of this
  node, and it is what \cref{lem:dc-ch8-2-1-jacobi-interior-data} was built for. It was
  \emph{not} the frames: an earlier estimate expected the cost to be in supplying frames and
  frame data, and that turned out to be the cheap part.

  \textbf{What this node does not discharge.} The single-chart clause is
  \cref{lem:dc-ch8-2-1-single-chart}, untouched. The no-conjugacy clause and the frame/$\phi$ data
  are hypotheses \emph{of this statement}, but they are no longer open: both are discharged one
  level up, by \cref{lem:dc-ch8-2-1-exp-norm-transfer-general-discharged}, using
  \cref{lem:dc-ch8-2-1-no-conjugate-normal} and \cref{lem:dc-ch8-2-1-cartan-frame-data}
  respectively. Finally this node is stated at one point $q=\exp_p(v)$; the passage to a
  neighborhood is \cref{lem:dc-ch8-2-1-local-isometry-packaging}.
\end{proof}

\begin{lemma}[an orthonormal basis for the metric at a point]
  \label{lem:dc-ch8-2-1-metric-ortho-basis}
  \lean{Riemannian.Jacobi.exists_metricOrthonormalBasis}
  \leanok
  Let $g$ be a Riemannian metric on $M$ and $q\in M$. Then $T_qM$ has an orthonormal basis for
  $\langle\cdot,\cdot\rangle_q$.

  The library previously had only chart-level producers, whose side conditions (membership in a
  chart target, or in a trivialization base set) exist solely to reach positive definiteness of
  the \emph{chart} Gram form. A Riemannian metric is positive definite at every point, so
  intrinsically the statement is unconditional in $q$.
\end{lemma}
\begin{proof}
  \leanok
  $\langle\cdot,\cdot\rangle_q$ is a symmetric bilinear form on the finite-dimensional space
  $T_qM$, positive definite by the definition of a Riemannian metric. Diagonalize it and
  normalize. $\square$
\end{proof}

\begin{lemma}[E. Cartan's $\phi_t$ as a construction, and the frames it matches]
  \label{lem:dc-ch8-2-1-cartan-frame-data}
  \lean{Riemannian.Jacobi.cartanTransportZero,
    Riemannian.Jacobi.cartanSeed,
    Riemannian.Jacobi.metricInner_cartanSeed,
    Riemannian.Jacobi.cartanPhi,
    Riemannian.Jacobi.cartanPhi_zero,
    Riemannian.Jacobi.cartanPhi_velocity,
    Riemannian.Jacobi.cartanPhi_isLinear,
    Riemannian.Jacobi.metricInner_cartanPhi,
    Riemannian.Jacobi.exists_cartanFrameData}
  \leanok
  \uses{lem:dc-ch8-2-1-phi,lem:dc-ch8-2-1-transported-frame,
    lem:dc-ch8-2-1-parallel-transport-map,lem:dc-ch8-2-1-metric-ortho-basis,
    lem:dc-ch8-2-1-velocity-frame}
  Let $M,\tilde M$ be complete, $p\in M$, $\tilde p\in\tilde M$, let $i:T_pM\to T_{\tilde p}\tilde
  M$ be a linear isometry, let $v\in T_pM$, and let $[a',b']$ be a window with $a'<0<1<b'$. Write
  $\gamma=\gamma_v$ and $\tilde\gamma=\gamma_{i(v)}$.

  Define $\phi^{\mathrm{Cartan}}_t=\tilde P_t\circ j\circ P_t^{-1}$ with
  $j=\tilde P_0^{-1}\circ i\circ P_0$ (the maps $P_0$, $j$ being
  \lean{Riemannian.Jacobi.cartanTransportZero} and \lean{Riemannian.Jacobi.cartanSeed}). Then for
  $t\in[a',b']$:
  \begin{itemize}
    \item $\phi^{\mathrm{Cartan}}_0=i$;
    \item $\phi^{\mathrm{Cartan}}_t(\gamma'(t))=\tilde\gamma'(t)$;
    \item $\phi^{\mathrm{Cartan}}_t$ is linear, and preserves the metric:
      $\langle\phi^{\mathrm{Cartan}}_tu,\phi^{\mathrm{Cartan}}_tw\rangle_{\tilde\gamma(t)}
      =\langle u,w\rangle_{\gamma(t)}$.
  \end{itemize}
  Moreover there exist parallel \emph{orthonormal} frames $e_1,\dots,e_n$ along $\gamma$ and
  $\tilde e_1,\dots,\tilde e_n$ along $\tilde\gamma$ on $[a',b']$ with
  $\phi^{\mathrm{Cartan}}_t(e_j(t))=\tilde e_j(t)$ for every $t\in[a',b']$.

  $\phi^{\mathrm{Cartan}}$ is a \emph{definition}, not an existential: its seed $j$ depends only on
  $g,\tilde g,p,\tilde p,v,i$ and the window, not on the frames. Only the frames are existential.
  This is load-bearing rather than cosmetic --- it is what lets
  \cref{lem:dc-ch8-2-1-exp-norm-transfer-general-discharged} state E. Cartan's curvature hypothesis
  \emph{about $\phi^{\mathrm{Cartan}}$} instead of quantifying over all admissible $\phi$; see the
  discussion there for why the quantified form is not merely inelegant but wrong.

  Together these supply every abstract-data hypothesis of
  \cref{lem:dc-ch8-2-1-exp-norm-transfer-general} \emph{except} the curvature match, which is
  E. Cartan's own hypothesis and must remain one.
\end{lemma}
\begin{proof}
  \leanok
  \uses{lem:dc-ch8-2-1-phi,lem:dc-ch8-2-1-transported-frame,
    lem:dc-ch8-2-1-metric-ortho-basis,lem:dc-ch8-2-1-velocity-frame}
  \cref{lem:dc-ch8-2-1-phi} builds $\tilde P_t\circ(-)\circ P_t^{-1}$ out of the parallel
  transports \emph{from the left endpoint} $a'$, and so seeding it with $i$ returns $\phi_{a'}=i$;
  but the transfer needs $\phi_0=i$, and $0$ is \emph{interior} to $[a',b']$.

  The fix does not require re-basing the transport. Conjugate the seed map instead: feed
  \cref{lem:dc-ch8-2-1-phi} not $i$ but
  $$
  j=\tilde P_0^{-1}\circ i\circ P_0 ,
  $$
  the map $i$ read from time $a'$ rather than from time $0$. Then
  $\phi_t=\tilde P_t\circ j\circ P_t^{-1}$, and at $t=0$ the outer transports cancel their own
  inverses, leaving $\phi_0=i$ exactly. That $j$ is again a linear isometry $T_{\gamma(a')}M\to
  T_{\tilde\gamma(a')}\tilde M$ follows because both transports preserve their Riemannian pairings
  (\cref{lem:dc-ch8-2-1-parallel-transport-map}) and $i$ is an isometry at $\gamma(0)=p$,
  $\tilde\gamma(0)=\tilde p$.

  With $j$ in hand, seed the frames by an orthonormal basis at $\gamma(a')$
  (\cref{lem:dc-ch8-2-1-metric-ortho-basis}) and apply
  \cref{lem:dc-ch8-2-1-transported-frame}; no basis needs transporting from $p$, since
  orthonormality is required at every time and parallel transport supplies it from any single one.
  Both remaining clauses are instances of the last claim of \cref{lem:dc-ch8-2-1-phi}: for
  $\phi_t(e_j(t))=\tilde e_j(t)$ apply it to the frames, whose seed identity at $a'$ the frame
  producer returns; for $\phi_t(\gamma'(t))=\tilde\gamma'(t)$ apply it to the velocity fields,
  which are parallel (\cref{lem:dc-ch8-2-1-velocity-frame}), the seed identity at $a'$ reducing
  through $P_0$, $\tilde P_0$ to $i(\gamma'(0))=i(v)=\tilde\gamma'(0)$. $\square$

  \textbf{Remark.} Unlike the Jacobi-side interior seeding
  (\cref{lem:dc-ch8-2-1-jacobi-interior-data}), no time reversal and no backward uniqueness are
  needed here: the inverted propagator already exists, parallel transport being invertible for the
  cheap reason that it is a metric isometry.
\end{proof}

\begin{lemma}[the transfer with its frames, $\phi$, and no-conjugacy discharged]
  \label{lem:dc-ch8-2-1-exp-norm-transfer-general-discharged}
  \lean{Riemannian.Jacobi.metricInner_mfderiv_eq_of_semiconjugacy_of_curvatureFormAt_of_isLocalDiffeomorphAt}
  \leanok
  \uses{lem:dc-ch8-2-1-exp-norm-transfer-general,lem:dc-ch8-2-1-cartan-frame-data,
    lem:dc-ch8-2-1-no-conjugate-normal}
  Same setting and same conclusion as \cref{lem:dc-ch8-2-1-exp-norm-transfer-general} --- $f$
  preserves the metric at $q=\exp_p(v)$ --- but with the frames, $\phi$, and \emph{both}
  no-conjugacy clauses discharged rather than assumed. In place of them it assumes only that
  $\exp_p$ and $\exp_{\tilde p}$ are local diffeomorphisms at $v$ and $i(v)$: the
  normal-neighborhood hypothesis. Two hypotheses survive: the single-chart clauses
  (\cref{lem:dc-ch8-2-1-single-chart}) and E. Cartan's curvature match.

  \textbf{The form of the curvature hypothesis.} The curvature match is stated about
  $\phi=\phi^{\mathrm{Cartan}}$ \emph{itself} --- the definition of
  \cref{lem:dc-ch8-2-1-cartan-frame-data} --- and so is do Carmo's own hypothesis and nothing more.

  This is worth stating because the obvious alternative is wrong. An earlier version of this node
  produced $\phi$ existentially and therefore had to require the match of \emph{every} $\phi$ that
  is a metric isometry along the geodesic pair, carries $\gamma'$ to $\tilde\gamma'$ and restricts
  to $i$ at $0$. That quantified hypothesis is strictly stronger, and not harmlessly so: at any
  interior $t\neq0$ one may replace $\phi_t$ by $\phi_t\circ\rho$ for any isometry $\rho$ of
  $T_{\gamma(t)}M$ fixing $\gamma'(t)$, and all three antecedents survive; comparing the two
  instances forces the curvature at $\gamma(t)$ to be invariant under all of $O(n-1)$ about the
  velocity. That isotropy holds in constant curvature and fails for a general $M$, so the
  quantified version had no variable-curvature content at all. Naming $\phi$ removes the hole.
\end{lemma}
\begin{proof}
  \leanok
  \uses{lem:dc-ch8-2-1-exp-norm-transfer-general,lem:dc-ch8-2-1-cartan-frame-data,
    lem:dc-ch8-2-1-no-conjugate-normal}
  Obtain the frames and $\phi$ from \cref{lem:dc-ch8-2-1-cartan-frame-data} and feed the curvature
  hypothesis the produced $\phi$, which satisfies its three antecedents by construction. Obtain
  the two no-conjugacy clauses from \cref{lem:dc-ch8-2-1-no-conjugate-normal} applied on each
  side. Then apply \cref{lem:dc-ch8-2-1-exp-norm-transfer-general}. $\square$
\end{proof}

\begin{lemma}[from pointwise metric preservation to a local isometry on $V$]
  \label{lem:dc-ch8-2-1-local-isometry-packaging}
  \uses{lem:dc-ch8-2-1-exp-norm-transfer-general-discharged,lem:dc-ch8-2-1-single-chart,
    lem:dc-ch8-2-1-exp-diff-zero,lem:dc-ch8-2-1-no-conjugate-zero}
  \notready
  \cref{lem:dc-ch8-2-1-exp-norm-transfer-general-discharged} preserves the metric at the
  \emph{single} point $q=\exp_p(v)$, for one $v$ whose geodesic satisfies the single-chart clause.
  \cref{thm:dc-ch8-2-1} asserts that $f$ is a local isometry on the whole normal neighborhood $V$.
  Bridging the two means running the pointwise statement for \emph{every} $w$ in a neighborhood of
  $0$ at once.

  \textbf{What this now costs.} Of the three clauses that used to have to be supplied uniformly in
  $w$, two are no longer obstacles. The frames and $\phi$ are produced along \emph{any} geodesic
  by \cref{lem:dc-ch8-2-1-cartan-frame-data}, and --- crucially --- the pointwise conclusion is
  pointwise in $q$, so an existential per $w$ suffices and no $w$-continuous frame field need be
  constructed. Non-conjugacy for each $\gamma_w$ now follows from
  \cref{lem:dc-ch8-2-1-no-conjugate-normal}, since $V$ is a \emph{normal} neighborhood and so
  $\exp_p$ is a local diffeomorphism at every $w$ of its model $U$; this is what an earlier draft
  of this node wrongly recorded as impossible in variable curvature.

  What remains is therefore the single-chart clause uniformly in $w$
  (\cref{lem:dc-ch8-2-1-single-chart}), plus the bookkeeping of assembling
  \lean{Riemannian.DCIsLocalIsometryAt} from the per-$w$ statements. In constant curvature this
  packaging is \cref{cor:dc-ch8-2-2}
  (\lean{Riemannian.Jacobi.dcIsLocalIsometryAt_of_semiconjugacy}), which escapes the single-chart
  difficulty because it only asserts that \emph{some} neighborhood works and so may shrink $V$
  into a chart; \cref{thm:dc-ch8-2-1} is handed $V$ and may not.

  A further gap, inherited from \cref{lem:dc-ch8-2-1-exp-norm-transfer-general-discharged}, is the
  form of the curvature hypothesis there: until $\phi$ is a definition rather than an existential,
  the assembled statement quantifies over all admissible $\phi$ and is weaker than do Carmo's.

  The base point $w=0$ is not part of this obligation: $df_p=i$ is already proved, curvature-free,
  by \cref{lem:dc-ch8-2-1-exp-diff-zero}.
\end{lemma}

\begin{lemma}[widening a single-chart time window]
  \label{lem:dc-ch8-2-1-window}
  \lean{Riemannian.Jacobi.exists_window_of_mem_chartAt_source}
  \leanok
  Let $\gamma:\mathbb R\to M$ be continuous and suppose $\gamma(t)$ lies in the source of a
  \emph{single} chart for every $t\in[0,1]$. Then it still does on some $[a',b']$ with $a'<0<1<b'$.

  Stated for an abstract continuous curve, so that one declaration serves both $\gamma$ and
  $\tilde\gamma$ of \cref{lem:dc-ch8-2-1-exp-norm-transfer-general}, which live over different
  manifolds.

  This lemma \emph{widens} a single-chart hypothesis; it does \emph{not} produce the chart. That
  such a chart exists at all is \cref{lem:dc-ch8-2-1-single-chart}, which remains open.
\end{lemma}
\begin{proof}
  \leanok
  The set $\{t:\gamma(t)\in\text{source}\}$ is open, being the preimage of an open set under a
  continuous map, and contains the compact $[0,1]$; so it contains a thickening
  $(-\delta,1+\delta)$ of it. Take $a'=-\delta/2$, $b'=1+\delta/2$. $\square$
\end{proof}

\begin{lemma}[a geodesic of a normal neighborhood, confined to one chart]
  \label{lem:dc-ch8-2-1-single-chart}
  \uses{lem:dc-ch8-2-1-frame-chart-interface,
    lem:dc-ch8-2-1-jacobi-norm-transfer-general,
    lem:dc-ch8-2-1-jacobi-norm-transfer-manifold}
  \notready
  \cref{lem:dc-ch8-2-1-frame-chart-interface} and
  \cref{lem:dc-ch8-2-1-jacobi-norm-transfer-general} both require that the geodesics
  $\gamma$, $\tilde\gamma$ each stay in the source of \emph{one} chart on the whole of
  $[a,b]$. For the geodesics $\gamma:[0,\ell]\to M$ of \cref{thm:dc-ch8-2-1}, which run
  across a normal neighborhood $V$ of $p$, this is not automatic: normal coordinates on $V$
  need not be a chart of the given atlas, and $\gamma([0,\ell])$ need not lie in a single
  chart source.

  What is required is therefore either a chart-chaining argument — cover the compact
  $[0,\ell]$ by finitely many parameter windows on each of which $\gamma$ stays in one chart
  source, run the comparison on each, and match at the junctions — or the hypothesis that
  $V$ lies in a chart source.

  One half of the second route is now available. \cref{lem:dc-ch8-2-1-window} widens a
  single-chart hypothesis on the compact $[0,1]$ to the two-sided window $[a',b']$ with
  $a'<0<1<b'$ that the transfer actually consumes, so under the hypothesis $V\subseteq$ a chart
  source the window costs nothing: $\gamma_w([0,1])\subseteq V$ by star-shapedness, and the
  widening supplies the rest. What is \emph{not} available, and is the whole of the present
  obligation, is the existence of that chart in the first place.

  The frames and $\phi_t$ do not obstruct the chaining: they are
  defined along all of $\gamma$ independently of any chart, and only their readings are
  chart-local. What must be threaded through the junctions is the ODE data, since
  \cref{lem:dc-ch8-2-1-jacobi-norm-transfer-general} propagates from initial data at $a$ and
  so cannot be localized time-by-time.

  This node and \cref{lem:dc-ch8-2-1-local-isometry-packaging} are the two obligations of
  \cref{thm:dc-ch8-2-1} that remain open; they are closely related, since the packaging node
  needs precisely this node's conclusion \emph{uniformly in $w$}, and closing this one for a
  single geodesic would not by itself close that one. The exponential-side chain is no longer
  among them: \cref{lem:dc-ch8-2-1-exp-norm-transfer-general} carries it from
  $|\tilde J|=|J|$ all the way to $\langle df_q u,df_q u'\rangle=\langle u,u'\rangle$, taking
  the present hypothesis as given.

  Note that \cref{cor:dc-ch8-2-2} does \emph{not} face the present obligation, because it
  asserts the mere existence of a neighborhood on which $f$ is an isometry and may shrink it
  into a chart; \cref{thm:dc-ch8-2-1} is handed $V$ and may not.
\end{lemma}

\begin{lemma}[the constant-curvature normal Jacobi field, any sign of $K_0$]
  \label{lem:dc-ch8-2-1-jacobi-const-solution}
  \lean{Riemannian.Jacobi.isJacobiFieldAlongOn_of_constantCurvature}
  \leanok
  \uses{def:dc-ch5-2-1}
  Let $M$ have constant sectional curvature $K_0$ (any sign), let $\gamma$ be a
  unit-speed geodesic, let $w$ be a \emph{parallel} field along $\gamma$ \emph{normal}
  to $\gamma'$, and let $h$ be any scalar solution of $h''+K_0 h=0$ (given by its
  derivative data $h'=Dh$, $Dh'=-K_0 h$). Then $J(t)=h(t)\,w(t)$ is a Jacobi field
  along $\gamma$, with covariant derivative $DJ(t)=Dh(t)\,w(t)$. A single statement thus
  covers the three curvature signs at once (do Carmo's
  $h=\sin(t\sqrt{K_0})/\sqrt{K_0}$, $h=t$, $h=\sinh(t\sqrt{-K_0})/\sqrt{-K_0}$).
\end{lemma}
\begin{proof}
  \leanok
  Since $w$ is parallel, $\tfrac{D}{dt}(h w)=h'\,w$ and
  $\tfrac{D}{dt}(Dh\,w)=Dh'\,w=-K_0 h\,w$; and by the constant-curvature model form
  $R(\gamma',w)\gamma'=K_0 w$ (do Carmo Ch.~4, Lemma 3.4, using $|\gamma'|^2=1$ and
  $\langle w,\gamma'\rangle=0$) together with homogeneity in the middle slot,
  $R(\gamma',h w)\gamma'=h K_0\,w$. Hence
  $\tfrac{D^2}{dt^2}(h w)+R(\gamma',h w)\gamma'=-K_0 h\,w+K_0 h\,w=0$: the pair
  $(h w, Dh\,w)$ solves the Jacobi system. This generalizes the special case
  $K_0>0$, $h=\sin(t\sqrt{K_0})/\sqrt{K_0}$ used in
  \cref{ex:dc-ch5-3-3} to an arbitrary scalar solution.
\end{proof}

\begin{lemma}[the norm of a constant-curvature Jacobi field is manifold-independent]
  \label{lem:dc-ch8-2-1-jacobi-norm-transfer}
  \lean{Riemannian.Jacobi.metricInner_jacobiField_eq_of_constantCurvature_of_speedSq,
    Riemannian.Jacobi.metricInner_jacobiField_transfer_of_constantCurvature_of_speedSq}
  \leanok
  \uses{lem:dc-ch8-2-1-jacobi-const-solution,cor:dc-ch5-2-5,def:dc-ch5-2-1}
  Let $\gamma:[0,\ell]\to M$ be a geodesic of squared speed $c=|\gamma'|^2\neq0$ on a
  manifold of constant sectional curvature $K_0$, let $J$ be the Jacobi field with $J(0)=0$
  and covariant derivative $DJ$, and let $h$ solve $h''+K_0c\,h=0$, $h(0)=0$, $h'(0)=1$.
  Writing $Z=DJ(0)$ and $a=\langle Z,\gamma'(0)\rangle$,
  $$
  |J(t)|^2 = \frac{a^2}{c}\,t^2 + h(t)^2\Big(|Z|^2-\frac{a^2}{c}\Big).
  $$
  The right-hand side involves \emph{no} manifold data beyond $K_0$ and $c$ (through $h$),
  the time $t$, and the two scalar invariants $a=\langle Z,\gamma'(0)\rangle$ and $|Z|^2$.
  Consequently, if $\tilde M$ also has constant curvature $K_0$, $\tilde\gamma$ is a geodesic
  of the \emph{same} squared speed $c$, $\tilde J$ is a Jacobi field with $\tilde J(0)=0$, and
  the initial invariants match ($\langle\tilde{DJ}(0),\tilde\gamma'(0)\rangle=a$,
  $|\tilde{DJ}(0)|^2=|Z|^2$), then $|\tilde J(t)|=|J(t)|$ for every $t\in[0,\ell]$. Since a
  linear isometry $i$ with $i(\gamma'(0))=\tilde\gamma'(0)$ carries $Z$ to $\tilde{DJ}(0)$
  and preserves both invariants --- and, being an isometry, automatically matches the two
  speeds --- this is precisely the norm-preserving (hence isometry) step of
  \cref{thm:dc-ch8-2-1} in the same-curvature case, and requires \emph{no}
  parallel-transport conjugation $\phi_t$. The unit-speed case $c=1$ recovers
  $|J(t)|^2=t^2a^2+h(t)^2(|Z|^2-a^2)$.
\end{lemma}
\begin{proof}
  \leanok
  Decompose $Z=\frac{a}{c}\,\gamma'(0)+Z_\perp$ with $Z_\perp=Z-\frac{a}{c}\gamma'(0)$; then
  $\langle Z_\perp,\gamma'(0)\rangle=a-\frac{a}{c}\cdot c=0$ and
  $|Z_\perp|^2=|Z|^2-\frac{a^2}{c}$. On a space of constant curvature $K_0$ a field normal to
  $\gamma'$ satisfies $R(\gamma',w)\gamma'=K_0|\gamma'|^2w=K_0c\,w$, so the normal Jacobi
  equation is the scalar equation $h''+K_0c\,h=0$. The tangential Jacobi field
  $\frac{a}{c}\,t\,\gamma'(t)$ (do Carmo Remark 2.2's $t\gamma'$) and the normal Jacobi field
  $h(t)\,w(t)$, with $w$ the parallel transport of $Z_\perp$
  (\cref{lem:dc-ch8-2-1-jacobi-const-solution}), both vanish at $t=0$ and have covariant
  derivatives $\frac{a}{c}\gamma'(0)$ resp. $Z_\perp$ there; their sum has initial data
  $(0,Z)$, so by uniqueness of Jacobi fields with prescribed initial data it equals $J$.
  Parallel transport is an isometry preserving orthogonality, so
  $\langle\gamma'(t),w(t)\rangle=\langle\gamma'(0),Z_\perp\rangle=0$, $|\gamma'(t)|^2=c$
  (geodesics have constant speed) and $|w(t)|^2=|Z_\perp|^2=|Z|^2-\frac{a^2}{c}$; expanding
  the squared norm of the orthogonal combination
  $\big|\frac{a}{c}\,t\,\gamma'(t)+h(t)\,w(t)\big|^2$ gives
  $\frac{a^2}{c^2}t^2\cdot c+h(t)^2(|Z|^2-\frac{a^2}{c})$, which is the claim. As this depends
  only on $K_0,c,t,a,|Z|^2$, equating it on the two sides yields $|\tilde J(t)|=|J(t)|$; by
  \cref{cor:dc-ch5-2-5} these are the values of $(d\exp_p)$ resp. $(d\exp_{\tilde p})$, so
  $df$ is norm-preserving.
\end{proof}

\begin{lemma}[the differential of $\exp$ is norm-preserving between spaces of the same
  constant curvature]
  \label{lem:dc-ch8-2-1-exp-norm-transfer}
  \lean{Riemannian.Jacobi.chartMetricInner_expDifferential_eq_metricInner_jacobiField,
    Riemannian.Jacobi.chartMetricInner_expDifferential_transfer_of_constantCurvature_of_speedSq}
  \leanok
  \uses{lem:dc-ch8-2-1-jacobi-norm-transfer,cor:dc-ch5-2-5,def:dc-ch5-2-1}
  Let $M$ and $\tilde M$ be complete manifolds of the same constant sectional curvature
  $K_0$, let $p\in M$, $\tilde p\in\tilde M$, and let $v\in T_p M$, $\tilde v\in T_{\tilde p}\tilde M$
  be \emph{nonzero} vectors of the same length, $|v|_p^2=|\tilde v|_{\tilde p}^2=c$. Let
  $Z\in T_p M$ and $\tilde Z\in T_{\tilde p}\tilde M$ have matching scalar invariants,
  $$
  \langle \tilde Z,\tilde v\rangle_{\tilde p}=\langle Z,v\rangle_p,
  \qquad |\tilde Z|^2_{\tilde p}=|Z|^2_p .
  $$
  Then the two exponential differentials have equal norms,
  $$
  \big|(d\exp_{\tilde p})_{\tilde v}(\tilde Z)\big| = \big|(d\exp_p)_v(Z)\big| .
  $$
  For $f=\exp_{\tilde p}\circ\,i\circ\exp_p^{-1}$ with $i$ a linear isometry and
  $\tilde v=i(v)$, $\tilde Z=i(Z)$, the invariants and the common length $c$ are preserved by
  $i$ automatically, so this is the mechanism behind the isometry step of
  \cref{thm:dc-ch8-2-1} in the same-curvature case, requiring \emph{no} parallel-transport
  conjugation $\phi_t$. Allowing every nonzero $v$, rather than only unit $v$, is what lets
  $q=\exp_p(v)$ range over a full normal neighborhood of $p$ rather than a single geodesic
  sphere --- the form \cref{cor:dc-ch8-2-2} requires.

  \emph{Scope of the \lean-tagged statement.} The formalized content is exactly the displayed
  norm equality. The passage from it to "$df_q$ is a linear isometry" --- assembling $f$ itself
  and applying the chain rule to $\exp_{\tilde p}\circ i\circ\exp_p^{-1}$ --- is not part of
  this node; see \cref{lem:dc-ch8-2-1-exp-local-diffeo} for the local inverse and
  \cref{lem:dc-ch8-2-1-polarization} for the passage from norms to inner products.
\end{lemma}
\begin{proof}
  \leanok
  By \cref{cor:dc-ch5-2-5} ($d(\exp_p)_v(Z)=Y_Z(1)$, read in a chart $\zeta$ around
  $\exp_p(v)$), each side is the value at time $1$ of the Jacobi field along $\gamma_v$
  resp.\ $\gamma_{\tilde v}$ with vanishing initial value and initial covariant derivative
  $Z$ resp.\ $\tilde Z$; a chart reading of a tangent vector is norm-faithful, so the
  chart-Gram norm of the differential is the intrinsic norm $|Y_Z(1)|$. Geodesics have
  constant speed (do Carmo Ch. 3, §2), and $\gamma_v$ has squared speed $|v|^2=c$, using that
  the intrinsic initial velocity of $\gamma_v$ is $v$ itself; so $\gamma_v$ and
  $\gamma_{\tilde v}$ have the same squared speed $c$, and $c\neq0$ because the metric is
  positive definite and $v\neq0$. \cref{lem:dc-ch8-2-1-jacobi-norm-transfer} then equates the
  two Jacobi norms, since the initial invariants match.
\end{proof}

\begin{lemma}[in constant curvature, no conjugate point where the scalar solution is nonzero]
  \label{lem:dc-ch8-2-1-no-conjugate}
  \lean{Riemannian.Jacobi.not_isConjugatePointAt_of_constantCurvature_of_speedSq,
    Riemannian.Jacobi.not_isConjugatePointAt_of_constantCurvature_of_lt_pi,
    Riemannian.Jacobi.exists_constCurvatureSol_ne_zero}
  \leanok
  \uses{lem:dc-ch8-2-1-jacobi-norm-transfer,lem:dc-ch8-2-1-jacobi-const-solution,def:dc-ch5-3-1}
  Let $\gamma:[0,\ell]\to M$ be a geodesic of squared speed $c\neq0$ on a manifold of constant
  sectional curvature $K_0$, and let $h$ solve $h''+K_0c\,h=0$, $h(0)=0$, $h'(0)=1$. If
  $h(\ell)\neq0$, then $\gamma(\ell)$ is \emph{not} conjugate to $\gamma(0)$ along $\gamma$.

  Only this direction is asserted. The converse --- that a zero of $h$ \emph{is} a conjugate
  point --- is not proved here; for $K_0>0$ it is witnessed at the first zero by
  \cref{ex:dc-ch5-3-3}, which exhibits $\gamma(\pi/\sqrt{K_0})$ as conjugate to $\gamma(0)$ in
  the unit-speed case (and additionally needs $n\geq2$).

  The criterion is sign-uniform. For $K_0\leq0$ the solution is $h(t)=t$ resp.
  $h(t)=\sinh(t\sqrt{-K_0c})/\sqrt{-K_0c}$, which never vanishes for $t>0$: there are no
  conjugate points at all, recovering the nonpositive-curvature statement of
  \cref{lem:dc-ch7-3-2}. For $K_0>0$ the solution is $h(t)=\sin(t\sqrt{K_0c})/\sqrt{K_0c}$,
  whose first positive zero is $\pi/\sqrt{K_0c}$, so the hypothesis holds throughout
  $0<\ell<\pi/\sqrt{K_0c}$ --- the interval that matters here, though $h$ is of course nonzero
  on further intervals beyond the first zero as well.

  Equivalently, in the $h$-free form: if
  $$
  K_0\,c\,\ell^2 < \pi^2,
  $$
  then $\gamma(\ell)$ is not conjugate to $\gamma(0)$. This single numerical condition covers
  all three signs at once --- it is vacuous for $K_0\leq0$, and for $K_0>0$ it reads
  $\ell\sqrt{K_0c}<\pi$.

  This is what supplies the no-conjugate-point hypothesis of
  \cref{lem:dc-ch8-2-1-exp-local-diffeo} on a whole \emph{ball} around $p$ in \emph{positive}
  constant curvature, where no nonpositive-curvature argument is available: applied along
  $\gamma_v$ at $\ell=1$ with $c=|v|^2$, the condition becomes $|v|<\pi/\sqrt{K_0}$. It is the
  input \cref{cor:dc-ch8-2-2} needs in order to reach every $q=\exp_p(v)$ of a normal
  neighborhood, and the radius is sharp by \cref{ex:dc-ch5-3-3}.
\end{lemma}
\begin{proof}
  \leanok
  Let $J$ be a Jacobi field along $\gamma$ with $J(0)=0$ and $J(\ell)=0$; we show $J\equiv0$,
  which contradicts the nontriviality required of a conjugacy witness. Write $Z=DJ(0)$ and
  $a=\langle Z,\gamma'(0)\rangle$. By \cref{lem:dc-ch8-2-1-jacobi-norm-transfer},
  $$
  0=|J(\ell)|^2=\frac{a^2}{c}\,\ell^2+h(\ell)^2\Big(|Z|^2-\frac{a^2}{c}\Big).
  $$
  Both summands are nonnegative: $c>0$ because $c=|\gamma'(0)|^2$ and the metric is positive
  definite with $c\neq0$; and $|Z|^2-\frac{a^2}{c}\geq0$ by Cauchy--Schwarz, being the squared
  norm of the component of $Z$ orthogonal to $\gamma'(0)$. A sum of two nonnegative reals
  vanishes only if both do. The first gives $\frac{a^2}{c}\ell^2=0$, hence $a=0$ since
  $\ell>0$ and $c>0$. The second then reads $h(\ell)^2|Z|^2=0$, and $h(\ell)\neq0$ forces
  $|Z|^2=0$, so $Z=DJ(0)=0$ by positive definiteness. Thus $J$ has vanishing initial data
  $(J(0),DJ(0))=(0,0)$, and by uniqueness of Jacobi fields with prescribed initial data
  $J\equiv0$ on $[0,\ell]$.

  For the $h$-free form, it remains to see that $K_0c\,\ell^2<\pi^2$ gives $h(\ell)\neq0$.
  Write $\kappa=K_0c$. If $\kappa<0$ then $h(t)=\sinh(t\sqrt{-\kappa})/\sqrt{-\kappa}$ and
  $\sinh$ is strictly increasing with $\sinh 0=0$, so $h(\ell)>0$; if $\kappa=0$ then
  $h(t)=t$ and $h(\ell)=\ell>0$; if $\kappa>0$ then $h(t)=\sin(t\sqrt\kappa)/\sqrt\kappa$, and
  $\kappa\ell^2<\pi^2$ with $\ell,\sqrt\kappa>0$ gives $0<\ell\sqrt\kappa<\pi$, whence
  $\sin(\ell\sqrt\kappa)>0$. $\square$
\end{proof}

\begin{lemma}[$\exp_p$ is a local diffeomorphism away from conjugate points]
  \label{lem:dc-ch8-2-1-exp-local-diffeo}
  \lean{Riemannian.Jacobi.isLocalDiffeomorphAt_expMapGlobal_of_not_conjugate}
  \leanok
  \uses{lem:dc-ch8-2-1-exp-equiv,prop:dc-ch5-3-5,lem:dc-ch7-3-2}
  Let $M$ be complete and $p\in M$. If the geodesic $\gamma_v$ from $p$ has no conjugate point
  of $p$ at parameter $1$, then $\exp_p$ is a $C^\infty$ local diffeomorphism at $v$; in
  particular it has a $C^\infty$ local inverse $\exp_p^{-1}$ defined near $\exp_p(v)$.

  This is what makes the map $f=\exp_{\tilde p}\circ\,i\circ\exp_p^{-1}$ of
  \cref{thm:dc-ch8-2-1} well defined and differentiable on a neighborhood of $p$, so that its
  differential $df_q$ is the honest one and the chain rule applies. No curvature hypothesis
  enters, which is essential for \cref{cor:dc-ch8-2-2}: it must also cover spaces of
  \emph{positive} constant curvature, where the nonpositive-curvature route of
  \cref{lem:dc-ch7-3-2} to the same conclusion is unavailable.
\end{lemma}
\begin{proof}
  \leanok
  By \cref{lem:dc-ch8-2-1-exp-equiv} the differential $d(\exp_p)_v$ is a continuous linear
  isomorphism: some chart $\zeta$ around $\exp_p(v)$ reads $\exp_p$ as a map with strict
  Fréchet derivative an isomorphism at $v$. That chart reading is $C^\infty$ on an open
  neighborhood of $v$, since $\exp_p$ is globally $C^\infty$ on a complete manifold, so the
  inverse function theorem makes the reading a local diffeomorphism at $v$. Composing with the
  chart inverse --- itself a local diffeomorphism --- recovers $\exp_p$ near $v$, which is
  therefore a local diffeomorphism at $v$ as well.
\end{proof}

\begin{lemma}[non-conjugacy from the normal neighborhood, in variable curvature]
  \label{lem:dc-ch8-2-1-no-conjugate-normal}
  \lean{Riemannian.Jacobi.not_isConjugatePointAt_globalGeodesic_of_injective_mfderiv,
    Riemannian.Jacobi.not_isConjugatePointAt_globalGeodesic_of_isLocalDiffeomorphAt,
    Riemannian.Jacobi.isLocalDiffeomorphAt_expMapGlobal_iff_not_isConjugatePointAt}
  \leanok
  \uses{lem:dc-ch8-2-1-exp-local-diffeo,lem:dc-ch8-2-1-mfderiv-bridge,prop:dc-ch5-3-5}
  Let $M$ be complete and $p\in M$, and let $v\in T_pM$. If $d(\exp_p)_v$ is injective --- in
  particular, if $\exp_p$ is a local diffeomorphism at $v$ --- then $1$ is \emph{not} a conjugate
  parameter along $\gamma_v$. No curvature hypothesis is needed.

  Together with \cref{lem:dc-ch8-2-1-exp-local-diffeo} this makes the two conditions
  \emph{equivalent}: $\exp_p$ is a $C^\infty$ local diffeomorphism at $v$ if and only if $1$ is
  not conjugate along $\gamma_v$.

  \textbf{Why this matters.} \cref{lem:dc-ch8-2-1-no-conjugate} produces non-conjugacy from a
  numerical condition and so is confined to constant curvature. The present lemma produces it
  instead from $v$ lying in the tangent-space model of a \emph{normal} neighborhood of $p$, and is
  therefore available in \emph{variable} curvature --- which is exactly the setting of
  \cref{thm:dc-ch8-2-1}, whose $V$ is normal by hypothesis. Earlier drafts of
  \cref{lem:dc-ch8-2-1-local-isometry-packaging} and of the discussion after
  \cref{lem:dc-ch8-2-1-exp-norm-transfer-general} recorded that no variable-curvature producer of
  the no-conjugacy clause existed; that is now false, and those passages have been corrected.
\end{lemma}
\begin{proof}
  \leanok
  \uses{lem:dc-ch8-2-1-mfderiv-bridge,prop:dc-ch5-3-5}
  By \cref{prop:dc-ch5-3-5} it suffices to show that the chart reading $D$ of $d(\exp_p)_v$ is
  injective. \cref{lem:dc-ch8-2-1-mfderiv-bridge} factors it as
  $$
  D = C_{\exp_p(v)\to\zeta}\circ d(\exp_p)_v ,
  $$
  and the coordinate change $C_{\exp_p(v)\to\zeta}$ is injective, its cocycle inverse being a left
  inverse; so $D$ is a composite of two injective maps.

  For the local-diffeomorphism form, the differential of a local diffeomorphism is a linear
  isomorphism, in particular injective. Note this is strictly cheaper than the special case
  \cref{lem:dc-ch8-2-1-no-conjugate-zero}, which additionally had to collapse $d(\exp_p)_0$ to the
  identity and transport along $\exp_p(0)=p$; here the differential stays abstract and only its
  injectivity is used. $\square$
\end{proof}

\begin{lemma}[polarization: a linear map preserving lengths preserves the metric]
  \label{lem:dc-ch8-2-1-polarization}
  \lean{Riemannian.Jacobi.metricInner_polarization,
    Riemannian.Jacobi.metricInner_transfer_of_norm_transfer}
  \leanok
  Let $q\in M$, $\tilde q\in\tilde M$, and let $L:T_qM\to T_{\tilde q}\tilde M$ be additive. If
  $L$ preserves squared lengths, $|L(w)|^2_{\tilde q}=|w|^2_q$ for every $w\in T_qM$, then it
  preserves the metric itself:
  $$
  \langle L(u),L(v)\rangle_{\tilde q} = \langle u,v\rangle_q
  \qquad\text{for all } u,v\in T_qM .
  $$
  This is the step from \cref{lem:dc-ch8-2-1-exp-norm-transfer}, which delivers only the
  equality of \emph{lengths} $|df_q(u)|=|u|$, to the statement that $f$ is a local isometry,
  which is phrased with the full inner product. Additivity cannot be dropped: a
  length-preserving map that is not additive need not preserve the metric. In the application
  $L=df_q$ is linear, so the hypothesis is automatic.
\end{lemma}
\begin{proof}
  \leanok
  A symmetric bilinear form over $\mathbb R$ is determined by its diagonal through the
  polarization identity: expanding $\langle u+v,u+v\rangle=\langle u,u\rangle+2\langle
  u,v\rangle+\langle v,v\rangle$ by bilinearity and symmetry gives
  $$
  \langle u,v\rangle=\tfrac12\big(\langle u+v,u+v\rangle-\langle u,u\rangle-\langle
  v,v\rangle\big).
  $$
  Apply this to $L(u),L(v)$ on $\tilde M$, use additivity to rewrite $L(u)+L(v)=L(u+v)$, then
  the length hypothesis at $u+v$, $u$ and $v$ to replace each of the three terms by its
  counterpart on $M$, and finally the same identity on $M$ read backwards:
  $$
  \langle L(u),L(v)\rangle=\tfrac12\big(|L(u+v)|^2-|L(u)|^2-|L(v)|^2\big)
  =\tfrac12\big(|u+v|^2-|u|^2-|v|^2\big)=\langle u,v\rangle. \qquad\square
  $$
\end{proof}

\begin{lemma}[$d(\exp_p)_v$ is an isomorphism, with the Jacobi identification]
  \label{lem:dc-ch8-2-1-exp-equiv}
  \lean{Riemannian.Jacobi.expDifferential_isEquiv_jacobi_of_not_conjugate}
  \leanok
  \uses{cor:dc-ch5-2-5,prop:dc-ch5-3-5}
  If the geodesic $\gamma_v$ from $p$ has no conjugate point of $p$ at parameter $1$, then
  $d(\exp_p)_v$ is a continuous linear isomorphism $T_pM\to T_{\exp_p(v)}M$, \emph{and} it still
  satisfies the identification of \cref{cor:dc-ch5-2-5}: it carries $\nabla J(0)$ to $J(1)$ for
  every Jacobi field $J$ along $\gamma_v$ with $J(0)=0$. Both facts come from the same strictly
  differentiable chart reading of $\exp_p$; the assembly of \cref{cor:dc-ch8-2-2} needs them
  simultaneously --- the isomorphism to invert $d\exp_p$ (so that every tangent vector at
  $q=\exp_p(v)$ is the endpoint value of a Jacobi field), and the identification to compute
  norms through \cref{lem:dc-ch8-2-1-exp-norm-transfer}.
\end{lemma}
\begin{proof}
  \leanok
  \cref{prop:dc-ch5-3-5} (\lstinline{expDifferential\_injective\_iff\_not\_conjugate}) makes the
  differential injective from the no-conjugate-point hypothesis, and an injective endomorphism
  of a finite-dimensional space is bijective, hence a continuous linear isomorphism. No
  curvature hypothesis enters: curvature is only ever used to \emph{produce} the
  no-conjugate-point hypothesis. The Jacobi clause is carried along unchanged from
  \cref{cor:dc-ch5-2-5}.
\end{proof}

\begin{lemma}[the chart reading of a differential is the coordinate change of the intrinsic one]
  \label{lem:dc-ch8-2-1-mfderiv-bridge}
  \lean{Riemannian.Jacobi.chartReading_fderiv_apply_eq,
    Riemannian.Jacobi.chartMetricInner_chartReading_fderiv_eq_metricInner_mfderiv,
    Riemannian.Jacobi.chartMetricInner_expDifferential_eq_metricInner_mfderiv}
  \leanok
  Let $F:T_pM\to M$ be $C^\infty$, let $\zeta\in M$ be a chart base point with $F(v)$ in the
  chart source, and suppose the chart-$\zeta$ reading $w\mapsto\varphi_\zeta(F(w))$ has
  differential $D$ at $v$. Then $D$ is the intrinsic differential of $F$ followed by the
  coordinate change of its own target fibre:
  $$
  D(Z)=C_{F(v)\to\zeta}\big(dF_v(Z)\big)\qquad\text{for all }Z ,
  $$
  and consequently, for the chart Gram form,
  $$
  \langle D(Z),D(Z')\rangle_{\zeta,\varphi_\zeta(F(v))}=\langle dF_v(Z),dF_v(Z')\rangle_{F(v)} .
  $$
  This is the bridge between the chart-Gram form in which \cref{lem:dc-ch8-2-1-exp-norm-transfer}
  is stated and the intrinsic form that \cref{cor:dc-ch8-2-2} requires; applied to $F=\exp_p$ it
  turns every chart-level statement about $d(\exp_p)_v$ into an intrinsic one.
\end{lemma}
\begin{proof}
  \leanok
  The chart map $\varphi_\zeta$ has an intrinsic differential at every point $y$ of its source,
  and that differential is exactly the tangent coordinate change $C_{y\to\zeta}$. Composing this
  with the intrinsic differential of $F$ at $v$ exhibits $C_{F(v)\to\zeta}\circ dF_v$ as a
  differential of the composite $\varphi_\zeta\circ F$ at $v$; since the source $T_pM$ is a vector
  space, that manifold differential is a Fréchet derivative of the chart reading. The Fréchet
  derivative is unique, so it coincides with $D$. The second display follows by applying the first
  to $Z$ and $Z'$ and using that a chart reading is norm-faithful: the chart Gram form at
  $\varphi_\zeta(x)$, evaluated on coordinate readings $C_{x\to\zeta}(w)$, returns
  $\langle w,w'\rangle_x$. $\square$
\end{proof}

\begin{lemma}[$d(\exp_p)_0$ is the identity]
  \label{lem:dc-ch8-2-1-exp-diff-zero}
  \lean{Riemannian.Jacobi.mfderiv_expMapGlobal_zero_apply}
  \leanok
  For $M$ complete and $p\in M$, the differential of $\exp_p$ at the origin of $T_pM$ is the
  identity: $d(\exp_p)_0(v)=v$ for every $v\in T_pM$.
\end{lemma}
\begin{proof}
  \leanok
  Fix $v$ and consider the ray $c:\mathbb R\to T_pM$, $c(t)=t\,v$, a linear map, so its own
  differential at $0$ is $c$ itself, and $c(1)=v$. Radial homogeneity of the exponential map gives
  $\exp_p(c(t))=\exp_p(t\,v)=\gamma_v(t)$, i.e. $\exp_p\circ\,c=\gamma_v$. Differentiating this
  identity at $t=0$ by the chain rule,
  $$
  d(\exp_p)_0\big(c(1)\big)=d(\gamma_v)_0(1)=\gamma_v'(0)=v ,
  $$
  the last equality because $\gamma_v$ is by construction the geodesic from $p$ with initial
  velocity $v$. Since $c(1)=v$, this reads $d(\exp_p)_0(v)=v$. $\square$
\end{proof}

\begin{lemma}[the base point is never conjugate to itself at parameter $1$]
  \label{lem:dc-ch8-2-1-no-conjugate-zero}
  \lean{Riemannian.Jacobi.not_isConjugatePointAt_globalGeodesic_zero}
  \leanok
  \uses{lem:dc-ch8-2-1-mfderiv-bridge,lem:dc-ch8-2-1-exp-diff-zero,prop:dc-ch5-3-5}
  For $M$ complete and $p\in M$, the parameter $1$ is not conjugate along the constant geodesic
  $\gamma_0$. No curvature hypothesis is needed. Consequently, by
  \cref{lem:dc-ch8-2-1-exp-local-diffeo}, $\exp_p$ is a $C^\infty$ local diffeomorphism at
  $0\in T_pM$ on \emph{any} complete Riemannian manifold --- the normal-neighborhood fact.

  This is what \cref{cor:dc-ch8-2-2} needs at the base point itself: the ball criterion
  \cref{lem:dc-ch8-2-1-no-conjugate} requires the speed $c=|v|^2$ to be nonzero and so says
  nothing at $v=0$.
\end{lemma}
\begin{proof}
  \leanok
  By \cref{prop:dc-ch5-3-5} the parameter $1$ fails to be conjugate along $\gamma_v$ exactly when
  the chart reading $D$ of $d(\exp_p)_v$ is injective, so it suffices to prove that for $v=0$. By
  \cref{lem:dc-ch8-2-1-mfderiv-bridge}, $D(Z)=C_{\exp_p(0)\to\zeta}\big(d(\exp_p)_0(Z)\big)$, and
  $\exp_p(0)=p$ while $d(\exp_p)_0$ is the identity by \cref{lem:dc-ch8-2-1-exp-diff-zero}. Hence
  $D=C_{p\to\zeta}$, which is injective, since the cocycle identity
  $C_{\zeta\to p}\circ C_{p\to\zeta}=C_{p\to p}=\mathrm{id}$ exhibits a left inverse. $\square$
\end{proof}

\begin{lemma}[non-conjugacy along $\gamma_v$, and surjectivity of $d(\exp_p)_v$]
  \label{lem:dc-ch8-2-1-no-conjugate-ball}
  \lean{Riemannian.Jacobi.not_isConjugatePointAt_globalGeodesic_of_constantCurvature_of_lt_pi,
    Riemannian.Jacobi.surjective_mfderiv_expMapGlobal_of_not_conjugate}
  \leanok
  \uses{lem:dc-ch8-2-1-no-conjugate,lem:dc-ch8-2-1-exp-local-diffeo}
  Let $M$ have constant curvature $K_0$ and let $v\in T_pM$ with $v\neq 0$. If
  $$
  K_0\,\langle v,v\rangle_p<\pi^2 ,
  $$
  then $1$ is not conjugate along $\gamma_v$. The condition is vacuous for $K_0\leq 0$ and reads
  $|v|<\pi/\sqrt{K_0}$ for $K_0>0$; it is sharp there, by \cref{ex:dc-ch5-3-3}. Moreover, whenever
  $1$ is not conjugate along $\gamma_v$, the differential $d(\exp_p)_v$ is surjective onto
  $T_{\exp_p(v)}M$.

  Surjectivity is the form in which do Carmo's step ``let $J$ be a Jacobi field along $\gamma$ with
  $J(0)=0$ and $J(\ell)=v$'' is used: every tangent vector at $q=\exp_p(v)$ is $d(\exp_p)_v(Z)$ for
  some $Z$, so a claim about $df_q$ need only be checked on such vectors.
\end{lemma}
\begin{proof}
  \leanok
  For the first claim, apply \cref{lem:dc-ch8-2-1-no-conjugate} along $\gamma_v$ at $\ell=1$: the
  geodesic $\gamma_v$ has constant speed $c=\langle v,v\rangle_p$, which is nonzero because the
  metric is positive definite and $v\neq 0$, and the numerical hypothesis
  $K_0\,c\,\ell^2<\pi^2$ is the assumption. For the second, \cref{lem:dc-ch8-2-1-exp-local-diffeo}
  makes $\exp_p$ a local diffeomorphism at $v$, and the differential of a local diffeomorphism is
  a linear isomorphism, in particular surjective. $\square$
\end{proof}

\begin{lemma}[the norm transfer, intrinsically]
  \label{lem:dc-ch8-2-1-exp-norm-transfer-mfderiv}
  \lean{Riemannian.Jacobi.metricInner_mfderiv_expMapGlobal_transfer_of_constantCurvature_of_speedSq}
  \leanok
  \uses{lem:dc-ch8-2-1-exp-norm-transfer,lem:dc-ch8-2-1-mfderiv-bridge,lem:dc-ch8-2-1-exp-equiv,
    lem:dc-ch8-2-1-jacobi-const-solution}
  Let $M$ and $\tilde M$ have the same constant curvature $K_0$, let $p\in M$, $\tilde p\in\tilde M$,
  and let $v,Z\in T_pM$, $\tilde v,\tilde Z\in T_{\tilde p}\tilde M$ with $v\neq 0$ and
  $$
  \langle v,v\rangle_p=\langle\tilde v,\tilde v\rangle_{\tilde p}=c,\qquad
  \langle \tilde Z,\tilde v\rangle_{\tilde p}=\langle Z,v\rangle_p,\qquad
  \langle \tilde Z,\tilde Z\rangle_{\tilde p}=\langle Z,Z\rangle_p .
  $$
  If $1$ is conjugate along neither $\gamma_v$ nor $\gamma_{\tilde v}$, then the two exponential
  differentials agree in length:
  $$
  \big|d(\exp_{\tilde p})_{\tilde v}(\tilde Z)\big|_{\exp_{\tilde p}(\tilde v)}
  =\big|d(\exp_p)_v(Z)\big|_{\exp_p(v)} .
  $$
  This is \cref{lem:dc-ch8-2-1-exp-norm-transfer} freed of its chart, and of the auxiliary scalar
  solution $h$.
\end{lemma}
\begin{proof}
  \leanok
  \cref{lem:dc-ch8-2-1-exp-equiv} supplies, on each side, a chart around the endpoint in which the
  reading of the exponential map is differentiable with derivative $D$ (resp.\ $\tilde D$) still
  carrying $\nabla J(0)$ to $J(1)$. \cref{lem:dc-ch8-2-1-exp-norm-transfer} then equates the two
  chart-Gram norms of $D(Z)$ and $\tilde D(\tilde Z)$, its scalar hypothesis being discharged by
  \cref{lem:dc-ch8-2-1-jacobi-const-solution}, which produces a solution of $h''+(K_0c)h=0$ with
  $h(0)=0$, $h'(0)=1$ for every $K_0$ and $c$. Finally
  \cref{lem:dc-ch8-2-1-mfderiv-bridge} rewrites each chart-Gram norm as the intrinsic norm of the
  corresponding intrinsic differential. $\square$
\end{proof}

\begin{lemma}[metric preservation along a semiconjugacy]
  \label{lem:dc-ch8-2-2-semiconjugacy}
  \lean{Riemannian.Jacobi.mfderiv_apply_eq_of_semiconjugacy_zero,
    Riemannian.Jacobi.metricInner_mfderiv_eq_of_semiconjugacy,
    Riemannian.Jacobi.dcIsLocalIsometryAt_of_semiconjugacy,
    Riemannian.Jacobi.dcIsLocalIsometryAt_of_semiconjugacy_self}
  \leanok
  \uses{lem:dc-ch8-2-1-exp-norm-transfer-mfderiv,lem:dc-ch8-2-1-no-conjugate-ball,
    lem:dc-ch8-2-1-no-conjugate-zero,lem:dc-ch8-2-1-exp-diff-zero,
    lem:dc-ch8-2-1-polarization,lem:dc-ch8-2-1-exp-local-diffeo}
  Let $M$ and $\tilde M$ have the same constant curvature $K_0$, let $i:T_pM\to T_{\tilde p}\tilde M$
  be a linear isometry, and let $W\subset T_pM$ be an open neighborhood of $0$ on which
  $K_0\langle w,w\rangle_p<\pi^2$. Suppose $f:M\to\tilde M$ is differentiable at each $\exp_p(w)$,
  $w\in W$, and \emph{semiconjugates} $\exp_p$ to $\exp_{\tilde p}$ through $i$:
  $$
  f\big(\exp_p(w)\big)=\exp_{\tilde p}\big(i(w)\big)\qquad\text{for all }w\in W .
  $$
  Then $df_p=i$, and $f$ is a local isometry at $p$: there is a neighborhood $V$ of $p$ with
  $\langle u,u'\rangle_q=\langle df_q(u),df_q(u')\rangle_{f(q)}$ for all $q\in V$ and
  $u,u'\in T_qM$.

  Taking the semiconjugacy as the hypothesis, rather than $f=\exp_{\tilde p}\circ\,i\circ\exp_p^{-1}$,
  is what makes the differential of $\exp_p^{-1}$ unnecessary: $d(\exp_p)_v$ is surjective, so $df_q$
  is pinned down on all of $T_qM$ by its values on the image of $d(\exp_p)_v$.
\end{lemma}
\begin{proof}
  \leanok
  Take $V=\exp_p(W')$, where $W'\subset W$ is the part of $W$ on which $\exp_p$ is inverted by the
  local inverse supplied by \cref{lem:dc-ch8-2-1-exp-local-diffeo} at $0$, available by
  \cref{lem:dc-ch8-2-1-no-conjugate-zero}. This $V$ is a neighborhood of $p$, and each $q\in V$ is
  $\exp_p(w)$ for a unique $w\in W'$.

  Suppose first $w\neq0$. Since $i$ is a linear isometry, $\langle i(w),i(w)\rangle_{\tilde p}
  =\langle w,w\rangle_p$, so the numerical hypothesis holds on both sides and by
  \cref{lem:dc-ch8-2-1-no-conjugate-ball} the parameter $1$ is conjugate along neither $\gamma_w$
  nor $\gamma_{i(w)}$; in particular $d(\exp_p)_w$ is surjective. Differentiating the semiconjugacy
  at $w$ by the chain rule --- both sides are differentiable, and the derivative of a map is unique
  --- gives
  $$
  df_q\big(d(\exp_p)_w(Z)\big)=d(\exp_{\tilde p})_{i(w)}\big(i(Z)\big)\qquad\text{for all }Z .
  $$
  The invariants $\langle i(Z),i(w)\rangle=\langle Z,w\rangle$ and $|i(Z)|^2=|Z|^2$ are preserved
  because $i$ is a linear isometry, so \cref{lem:dc-ch8-2-1-exp-norm-transfer-mfderiv} equates the
  lengths of the two sides. As $Z$ ranges over $T_pM$ the vector $d(\exp_p)_w(Z)$ ranges over all of
  $T_qM$, so $df_q$ preserves lengths; being linear, it preserves the metric, by
  \cref{lem:dc-ch8-2-1-polarization}.

  If instead $w=0$ then $q=\exp_p(0)=p$. The same chain rule at $0$, together with
  $d(\exp_p)_0=d(\exp_{\tilde p})_0=\mathrm{id}$ (\cref{lem:dc-ch8-2-1-exp-diff-zero}) and
  $i(0)=0$, gives $df_p=i$; and $f(p)=\exp_{\tilde p}(i(0))=\tilde p$. The claim at $q=p$ is then
  exactly the hypothesis that $i$ is a linear isometry. $\square$
\end{proof}

\begin{theorem}[E. Cartan]
  \label{thm:dc-ch8-2-1}
  \dcref{ch8:2.1}
  \uses{cor:dc-ch5-2-5,lem:dc-ch8-2-1-jacobi-transfer,lem:dc-ch8-2-1-parallel-ortho-frame,
    lem:dc-ch8-2-1-transported-frame,lem:dc-ch8-2-1-hmatch-transfer,
    lem:dc-ch8-2-1-curvature-bridge,lem:dc-ch8-2-1-frame-coeff-bridge,
    lem:dc-ch8-2-1-jacobi-norm-transfer-general,lem:dc-ch8-2-1-frame-chart-interface,
    lem:dc-ch8-2-1-jacobi-norm-transfer-manifold,lem:dc-ch8-2-1-single-chart,
    lem:dc-ch8-2-1-exp-norm-transfer-general,lem:dc-ch8-2-1-phi,
    lem:dc-ch8-2-1-parallel-transport-map,lem:dc-ch8-2-1-exp-norm-transfer,
    lem:dc-ch8-2-1-exp-equiv,lem:dc-ch8-2-1-exp-local-diffeo,
    lem:dc-ch8-2-1-polarization,lem:dc-ch8-2-1-no-conjugate,
    lem:dc-ch8-2-1-mfderiv-bridge,lem:dc-ch8-2-1-exp-diff-zero,
    lem:dc-ch8-2-1-no-conjugate-zero,lem:dc-ch8-2-1-no-conjugate-ball,
    lem:dc-ch8-2-1-exp-norm-transfer-mfderiv,
    lem:dc-ch8-2-1-jacobi-reversal,lem:dc-ch8-2-1-jacobi-backward-uniqueness,
    lem:dc-ch8-2-1-jacobi-interior-data,lem:dc-ch8-2-1-local-isometry-packaging}
  \notready
  With the notation above, if for all $q\in V$ and all $x,y,u,v\in T_q(M)$ we have

  $$
  \langle R(x,y)u,v\rangle=\big\langle\tilde{R}(\phi_t(x),\phi_t(y))\phi_t(u),\phi_t(v)\big\rangle,
  $$

  then $f:V\to f(V)\subset\tilde{M}$ is a local isometry and $df_p=i$.
\end{theorem}
\begin{proof}
  Let $q\in V$ and let $\gamma:[0,\ell]\to M$ be a normalized geodesic with
  $\gamma(0)=p$, $\gamma(\ell)=q$. Let $v\in T_q(M)$ and let $J$ be a Jacobi field
  along $\gamma$ with $J(0)=0$ and $J(\ell)=v$. Let $e_1,\dots,e_n=\gamma'(0)$ be an
  orthonormal basis of $T_p(M)$ and let $e_i(t)$ be the parallel transport of $e_i$
  along $\gamma$. Writing $J(t)=\sum_i y_i(t)e_i(t)$, the Jacobi equation gives

  $$
  y_j''+\sum_i\langle R(e_n,e_i)e_n,e_j\rangle y_i=0,\qquad j=1,\dots,n.
  $$

  Let $\tilde{\gamma}:[0,\ell]\to\tilde{M}$ be the normalized geodesic with
  $\tilde{\gamma}(0)=\tilde{p}$, $\tilde{\gamma}'(0)=i(\gamma'(0))$, and set
  $\tilde{J}(t)=\phi_t(J(t))$, $\tilde{e}_j(t)=\phi_t(e_j(t))$. By linearity of
  $\phi_t$, $\tilde{J}(t)=\sum_i y_i(t)\tilde{e}_i(t)$. Since by hypothesis
  $\langle R(e_n,e_i)e_n,e_j\rangle=\langle\tilde{R}(\tilde{e}_n,\tilde{e}_i)\tilde{e}_n,\tilde{e}_j\rangle$,
  $\tilde{J}$ satisfies the same equation, hence $\tilde{J}$ is a Jacobi field along
  $\tilde{\gamma}$ with $\tilde{J}(0)=0$. Since parallel transport is an isometry,
  $|\tilde{J}(\ell)|=|J(\ell)|$. It remains to show that $\tilde{J}(\ell)=df_q(v)=df_q(J(\ell))$.
  Since $\tilde{J}(t)=\phi_t(J(t))$, $\tilde{J}'(0)=i(J'(0))$. Since $J$ and
  $\tilde{J}$ are Jacobi fields vanishing at $t=0$, by \cref{cor:dc-ch5-2-5},

  $$
  J(t)=(d\exp_p)_{t\gamma'(0)}(tJ'(0)),\qquad\tilde{J}(t)=(d\exp_{\tilde{p}})_{t\tilde{\gamma}'(0)}(t\tilde{J}'(0)).
  $$

  Therefore

  $$
  \tilde{J}(\ell)=(d\exp_{\tilde{p}})_{\ell\tilde{\gamma}'(0)}\circ i\circ\big((d\exp_p)_{\ell\gamma'(0)}\big)^{-1}(J(\ell))=df_q(J(\ell)),
  $$

  which proves the assertion. $\square$

  Observe that the same proof shows that if $\exp_p$ and $\exp_{\tilde{p}}$ are
  diffeomorphisms, then, under the conditions of Theorem 2.1, $f$ is defined on all
  of $M$ and is an isometry. The theorem implies that the metric is, in a certain
  sense, determined locally by the curvature. An equivalent assertion was made by
  Riemann; the first proof of the local theorem was presented by E. Cartan. A global
  version was given by W. Ambrose in 1956.
\end{proof}

\begin{corollary}[Corollary 2.2]
  \label{cor:dc-ch8-2-2}
  \dcref{ch8:2.2}
  \lean{Riemannian.Jacobi.exists_dcIsLocalIsometryAt_of_constantCurvature}
  \leanok
  \uses{lem:dc-ch8-2-1-jacobi-transfer-const,lem:dc-ch8-2-1-velocity-frame,
    lem:dc-ch8-2-1-jacobi-norm-transfer,lem:dc-ch8-2-1-jacobi-const-solution,
    lem:dc-ch8-2-1-exp-norm-transfer,lem:dc-ch8-2-1-exp-equiv,
    lem:dc-ch8-2-1-exp-local-diffeo,lem:dc-ch8-2-1-polarization,
    lem:dc-ch8-2-1-no-conjugate,lem:dc-ch8-2-1-mfderiv-bridge,
    lem:dc-ch8-2-1-exp-diff-zero,lem:dc-ch8-2-1-no-conjugate-zero,
    lem:dc-ch8-2-1-no-conjugate-ball,lem:dc-ch8-2-1-exp-norm-transfer-mfderiv,
    lem:dc-ch8-2-2-semiconjugacy}
  Let $M$ and $\tilde{M}$ be spaces with the same constant curvature and the same
  dimension $n$. Let $p\in M$ and $\tilde{p}\in\tilde{M}$. Choose arbitrary
  orthonormal bases $\{e_j\}$ of $T_p(M)$ and $\{\tilde{e}_j\}$ of $T_{\tilde{p}}(\tilde{M})$,
  $j=1,\dots,n$. Then there exist a neighborhood $V\subset M$ of $p$, a neighborhood
  $\tilde{V}\subset\tilde{M}$ of $\tilde{p}$, and an isometry $f:V\to\tilde{V}$ such
  that $df_p(e_j)=\tilde{e}_j$.
\end{corollary}
\begin{proof}
  Choose the isometry $i$ of the theorem in such a way that
  $i(e_j)=\tilde{e}_j$. The condition on the curvature is immediately verified, and
  the conclusion follows from \cref{thm:dc-ch8-2-1}. $\square$

  \leanok
  Since $M$ and $\tilde M$ share the constant curvature $K_0$, the curvature-matching hypothesis
  of \cref{thm:dc-ch8-2-1} is automatic and the parallel-transport conjugation $\phi_t$ is not
  needed, so the corollary does not in fact wait on \cref{thm:dc-ch8-2-1} in its full generality.
  Concretely: let $\exp_p^{-1}$ be the local inverse of \cref{lem:dc-ch8-2-1-exp-local-diffeo} at
  $0$, which exists with no curvature hypothesis by \cref{lem:dc-ch8-2-1-no-conjugate-zero}, and
  set
  $$
  f=\exp_{\tilde p}\circ\,i\circ\exp_p^{-1} .
  $$
  On the window $W=\operatorname{dom}(\exp_p^{-1})\cap\{w: K_0\langle w,w\rangle_p<\pi^2\}$ this
  $f$ satisfies $f\circ\exp_p=\exp_{\tilde p}\circ\,i$ by construction. The window is an open
  neighborhood of $0$: the second factor is the preimage of $(-\infty,\pi^2)$ under
  $w\mapsto K_0\langle w,w\rangle_p$, which is continuous because the metric at a point is a
  continuous bilinear form; it is all of $T_pM$ when $K_0\leq0$, and the ball
  $|w|<\pi/\sqrt{K_0}$ when $K_0>0$, where it is sharp by \cref{ex:dc-ch5-3-3}. Also $f$ is
  differentiable on $\exp_p(W)$, being a composite of $\exp_p^{-1}$ (smooth on its domain), the
  linear map $i$, and $\exp_{\tilde p}$ (smooth). So
  \cref{lem:dc-ch8-2-2-semiconjugacy} applies and gives both that $f$ is a local isometry at $p$
  and that $df_p=i$; choosing $i(e_j)=\tilde e_j$ turns the latter into $df_p(e_j)=\tilde e_j$.
  (Such an $i$ exists: both bases are orthonormal for their respective metrics and the two spaces
  have the same dimension, so the linear map carrying $e_j$ to $\tilde e_j$ intertwines the two
  metrics, by bilinearity. The Lean statement is the correspondingly \emph{stronger} ``for every
  metric-intertwining linear isomorphism $i$'' form; the passage from a pair of orthonormal bases
  to such an $i$ is elementary linear algebra and is not separately formalized --- tracked as
  I-0320.)
  Finally $f$ is a local diffeomorphism at $p$, being a composite of three local
  diffeomorphisms --- $\exp_p^{-1}$ at $p$, the linear isomorphism $i$, and $\exp_{\tilde p}$ at
  $0$ (again \cref{lem:dc-ch8-2-1-no-conjugate-zero}, now on $\tilde M$) --- so it restricts to a
  diffeomorphism of a neighborhood $V$ of $p$ onto a neighborhood $\tilde V$ of
  $f(p)=\tilde p$. $\square$
\end{proof}

\begin{corollary}[Corollary 2.3]
  \label{cor:dc-ch8-2-3}
  \dcref{ch8:2.3}
  \lean{Riemannian.Jacobi.exists_dcIsLocalIsometryAt_of_constantCurvature_self}
  \leanok
  \uses{lem:dc-ch8-2-1-jacobi-transfer-const,lem:dc-ch8-2-1-velocity-frame,
    lem:dc-ch8-2-1-jacobi-norm-transfer,lem:dc-ch8-2-1-jacobi-const-solution,
    lem:dc-ch8-2-1-exp-norm-transfer,lem:dc-ch8-2-1-exp-equiv,
    lem:dc-ch8-2-1-exp-local-diffeo,lem:dc-ch8-2-1-polarization,
    lem:dc-ch8-2-1-no-conjugate,lem:dc-ch8-2-1-mfderiv-bridge,
    lem:dc-ch8-2-1-exp-diff-zero,lem:dc-ch8-2-1-no-conjugate-zero,
    lem:dc-ch8-2-1-no-conjugate-ball,lem:dc-ch8-2-1-exp-norm-transfer-mfderiv,
    lem:dc-ch8-2-2-semiconjugacy}
  Let $M$ be a space of constant curvature and let $p$ and $q$ be any two points of
  $M$. Let $\{e_j\}$ and $\{f_j\}$ be arbitrary orthonormal bases of $T_p(M)$ and
  $T_q(M)$, respectively. Then there exist neighborhoods $U$ of $p$ and $V$ of $q$,
  and an isometry $g:U\to V$ such that $dg_p(e_j)=f_j$.

  Related to the above is the problem of deciding whether a diffeomorphism
  $f:M\to\tilde{M}$ which preserves the curvature in the sense that

  $$
  \langle R(X,Y)Z,T\rangle_p=\big\langle\tilde{R}(df_p(X),df_p(Y))df_p(Z),df_p(T)\big\rangle_{f(p)}
  $$

  for all $p\in M$ and all $X,Y,Z,T\in T_p(M)$ is an isometry. In dimension two this
  is a kind of converse to the Theorem Egregium of Gauss, and is false, even in the
  compact case (Fig. 1). For $n=\dim M=\dim\tilde{M}\geq 4$, the problem
  surprisingly admits an affirmative solution (Kulkarni); for $n=3$, the problem was
  treated by Yau.
\end{corollary}
\begin{proof}
  \leanok
  \uses{cor:dc-ch8-2-2}
  Apply \cref{cor:dc-ch8-2-2} with $\tilde M=M$ and $\tilde p=q$, taking for $i$ the linear
  isometry $T_pM\to T_qM$ carrying $e_j$ to $f_j$; such an $i$ exists because both bases are
  orthonormal (elementary linear algebra, not separately formalized --- see I-0320; the Lean
  statement quantifies over every metric-intertwining $i$, hence is stronger). $\square$
\end{proof}

\section{Hyperbolic space}

We give now an example of a space of constant curvature $-1$. Consider the
half-space of $\mathbb{R}^n$ given by

$$
H^n=\{(x_1,\dots,x_n)\in\mathbb{R}^n;\ x_n>0\}
$$

and introduce on $H^n$ the metric

$$
g_{ij}(x_1,\dots,x_n)=\frac{\delta_{ij}}{x_n^2}.\qquad(1)
$$

It is clear that $H^n$ is simply connected. We are going to show that $H^n$ with
the metric $(1)$ has constant sectional curvature equal to $-1$, and that $H^n$ is
complete. $H^n$ is called the \emph{hyperbolic space} of dimension $n$.

\begin{definition}[hyperbolic space $H^n$ as a Riemannian manifold]
  \label{def:dc-ch8-3-hyperbolic-space}
  \lean{Riemannian.opensConformalMetric,
    Riemannian.Hyperbolic.hyperbolicMetric,
    Riemannian.Hyperbolic.hyperbolicMetric_apply,
    Riemannian.Hyperbolic.hyperbolicMetric_gram}
  \leanok
  The \emph{hyperbolic space} $H^n$ is the open half-space
  $H^n=\{x\in\mathbb R^n : x_n>0\}$ carrying the metric $g_{ij}=\delta_{ij}/x_n^2$
  of $(1)$, that is, the conformal rescaling of the ambient Euclidean metric by
  the positive smooth factor $1/x_n^2$: for $u,v\in T_pH^n$,
  $\langle u,v\rangle_g=\tfrac{1}{x_n^2}\langle u,v\rangle$. This is a Riemannian
  manifold. More generally, given an open subset of a (complete) real
  inner-product space and a positive smooth function $\varphi$ on it, the
  pointwise rescaling $\langle u,v\rangle\mapsto\varphi\,\langle u,v\rangle$ of
  the ambient Euclidean metric is a Riemannian metric — the conformal metric with
  conformal factor $\varphi$ (positive-definiteness and von-Neumann boundedness
  are inherited from the flat metric, whose metric ball is a Euclidean ball, and
  smoothness of the metric section from that of $\varphi$).
\end{definition}

A good part of the calculation can be carried out in a more general situation. Two
metrics $\langle\,,\rangle$ and $\langle\!\langle\,,\rangle\!\rangle$ on a
differentiable $M$ are \emph{conformal} if there exists a positive differentiable
function $f:M\to\mathbb{R}$ such that
$\langle u,v\rangle_p=f(p)\langle\!\langle u,v\rangle\!\rangle_p$ for all $p\in M$
and $u,v\in T_p(M)$. The metric $(1)$ of $H^n$ is conformal to the usual metric of
Euclidean space. Consider on $H^n$ the metric $g_{ij}=\delta_{ij}/F^2$, where $F$
is a positive differentiable function on $H^n$; write $g^{ij}=F^2\delta_{ij}$ for
the inverse matrix, and put $\log F=f$, $\partial f/\partial x_j=f_j$. Then

$$
\frac{\partial g_{ik}}{\partial x_j}=-2\frac{\delta_{ik}}{F^2}f_j.
$$

A computation of the Christoffel symbols gives

$$
\Gamma_{ij}^k=-\delta_{jk}f_i-\delta_{ki}f_j+\delta_{ij}f_k;
$$

hence if all three indices are distinct, $\Gamma_{ij}^k=0$, while if two indices
are equal,

$$
\Gamma_{ij}^i=-f_j,\quad\Gamma_{ii}^j=f_j,\quad\Gamma_{ij}^j=-f_i,\quad\Gamma_{ii}^i=-f_i.
$$

A direct calculation of the curvature coefficients yields

$$
F^2 R_{ijij}=-\sum_\ell f_\ell^2+f_i^2+f_j^2+f_{ii}+f_{jj},
$$

and $R_{ijk\ell}=0$ if all four indices are distinct; if any three indices are
distinct,

$$
R_{ijk}^i=-f_k f_j-f_{kj},\quad R_{ijk}^j=f_i f_k+f_{ki},\quad R_{ijk}^k=0.\qquad(2)
$$

The sectional curvature with respect to the plane generated by
$\partial/\partial x_i,\partial/\partial x_j$ (which are orthogonal) is

$$
K_{ij}=\frac{R_{ijij}}{g_{ii}g_{jj}}=R_{ijij}F^4=\Big(-\sum_\ell f_\ell^2+f_i^2+f_j^2+f_{ii}+f_{jj}\Big)F^2.
$$

Specializing to $F^2=x_n^2$, so that $f=\log x_n$: if $i\neq n$ and $j\neq n$,
$K_{ij}=-\tfrac{1}{x_n^2}x_n^2=-1$; if $i=n,j\neq n$,
$K_{nj}=(-f_n^2+f_n^2+f_{nn})F^2=-\tfrac{1}{x_n^2}x_n^2=-1$; and similarly if
$i\neq n,j=n$, $K_{in}=-1$. Using $(2)$ and \cref{cor:dc-ch4-3-5}, we conclude
that the sectional curvature of $H^n$ is a constant equal to $-1$.

\begin{lemma}[Christoffel symbols and coordinate curvature of a conformal metric]
  \label{lem:dc-ch8-3-conformal-curvature}
  \lean{Riemannian.Hyperbolic.christoffelFromMetric_eq,
    Riemannian.Hyperbolic.Rcoeff_diag,
    Riemannian.Hyperbolic.hyperbolic_sectionalCurvature}
  \leanok
  For the conformal metric $g_{ij}=\delta_{ij}/F^2$ ($f=\log F$), the Christoffel
  symbols computed from the Koszul formula
  $\Gamma^k_{ij}=\tfrac12\sum_\ell g^{k\ell}(\partial_i g_{j\ell}+\partial_j g_{i\ell}-\partial_\ell g_{ij})$
  equal $\Gamma^k_{ij}=-\delta_{jk}f_i-\delta_{ki}f_j+\delta_{ij}f_k$, and the
  associated coordinate curvature coefficients satisfy, for $i\neq j$,
  $$F^2R_{ijij}=-\sum_\ell f_\ell^2+f_i^2+f_j^2+f_{ii}+f_{jj}.$$
  Specializing to $F=x_n$ (so $f=\log x_n$, $f_a=\delta_{an}/x_n$,
  $f_{ab}=-\delta_{an}\delta_{bn}/x_n^2$) gives, for every \emph{coordinate}
  $2$-plane spanned by $\partial/\partial x_i,\partial/\partial x_j$ with $i\neq j$,
  the coordinate-frame sectional curvature $K_{ij}=R_{ijij}F^4=-1$. (This is the
  coordinate identity; extending it to all $2$-planes and to the intrinsic
  curvature of the manifold is \cref{prop:dc-ch8-3-const-curv}.)
\end{lemma}
\begin{proof}
  \leanok
  The Koszul identity is a direct computation from
  $g^{k\ell}=\delta_{k\ell}F^2$ and $\partial_m g_{ab}=-2\delta_{ab}f_m/F^2$.
  The curvature coefficient is
  $R^s_{ijk}=\partial_j\Gamma^s_{ik}-\partial_i\Gamma^s_{jk}+\sum_\ell(\Gamma^\ell_{ik}\Gamma^s_{j\ell}-\Gamma^\ell_{jk}\Gamma^s_{i\ell})$;
  substituting the closed form for $\Gamma$ and contracting the Kronecker deltas
  yields $R^j_{iji}=-\sum_\ell f_\ell^2+f_i^2+f_j^2+f_{ii}+f_{jj}$, which is
  $F^2R_{ijij}$ since $g_{jj}=1/F^2$. For $F=x_n$ one has
  $(f_a)^2+f_{aa}=0$ at every index, so the bracketed sum collapses to
  $-\sum_\ell f_\ell^2=-1/x_n^2$, and $K_{ij}=R_{ijij}F^4=(F^2R_{ijij})F^2=-1$.
\end{proof}

\begin{lemma}[the Riemann curvature tensor in chart coordinates]
  \label{lem:dc-ch8-3-chart-curvature}
  \lean{Riemannian.leviCivita_curvature_chartFrame_expansion}
  \leanok
  Let $g$ be a Riemannian metric on $M$ and $\nabla$ its Levi-Civita connection.
  In the coordinate frame $\partial_i$ of a chart at a point $p$, the intrinsic
  curvature operator is computed from the chart Christoffel symbols
  $\Gamma^k_{ij}$ by the classical coordinate formula
  $$
  R(\partial_i,\partial_j)\partial_k
    = \sum_\ell\Big(\partial_j\Gamma^\ell_{ik}-\partial_i\Gamma^\ell_{jk}
      +\sum_s(\Gamma^s_{ik}\Gamma^\ell_{js}-\Gamma^s_{jk}\Gamma^\ell_{is})\Big)\partial_\ell,
  $$
  where $\partial_j\Gamma^\ell_{ik}$ is the chart partial derivative of the
  coordinate Christoffel function. With do Carmo's sign convention
  $R(X,Y)Z=\nabla_Y\nabla_X Z-\nabla_X\nabla_Y Z+\nabla_{[X,Y]}Z$ the bracketed
  coefficient is exactly the $R^\ell_{ijk}$ of $(2)$.
\end{lemma}
\begin{proof}
  \leanok
  In a coordinate chart the frame vectors have vanishing pairwise Lie brackets,
  $[\partial_i,\partial_j]=0$, so the bracket term
  $\nabla_{[\partial_i,\partial_j]}\partial_k$ drops out. The field
  $\nabla_{\partial_j}\partial_k$ is, throughout the chart, the frame combination
  $\sum_m\Gamma^m_{jk}\partial_m$ (formula $(10)$ of Chap.~2, the Koszul collapse
  on the coordinate frame, valid at every chart point). Differentiating this
  along $\partial_i$ by the Leibniz rule for $\nabla$ gives
  $\nabla_{\partial_i}\nabla_{\partial_j}\partial_k
   =\sum_\ell\big(\partial_i\Gamma^\ell_{jk}
   +\sum_s\Gamma^s_{jk}\Gamma^\ell_{is}\big)\partial_\ell$, using
  $\nabla_{\partial_i}\partial_m=\sum_\ell\Gamma^\ell_{im}\partial_\ell$ for the
  inner derivative. Subtracting the term with $i$ and $j$ interchanged yields the
  stated coefficients.
\end{proof}

\begin{lemma}[hyperbolic chart Gram and inverse-Gram in the abstract chart frame]
  \label{lem:dc-ch8-3-hyperbolic-chart-gram}
  \lean{Riemannian.Hyperbolic.hyperbolic_chartGramMatrix,
    Riemannian.Hyperbolic.hyperbolic_chartGramMatrix_eq,
    Riemannian.Hyperbolic.hyperbolic_chartInvGramMatrix,
    Riemannian.Hyperbolic.finBasisGramMatrix_det_isUnit}
  \leanok
  \uses{def:dc-ch8-3-hyperbolic-space}
  Because $H^n=\{x_n>0\}$ is an open subset of $\mathbb R^n$, its tangent bundle is
  canonically trivial and the abstract chart frame vector
  $\partial_i=$ \emph{finBasis}$_i$ (the frame underlying \cref{lem:dc-ch8-3-chart-curvature})
  is literally the constant basis vector. Hence the chart Gram matrix of the
  hyperbolic metric is the conformal rescaling
  $G_{ij}(x)=x_n^{-2}\,B_{ij}$ of the constant Gram matrix
  $B_{ij}=\langle\text{finBasis}_i,\text{finBasis}_j\rangle$ of the abstract basis,
  and its inverse is $G^{ij}(x)=x_n^{2}\,B^{ij}$ (with $B$ invertible, being the
  Gram matrix of a basis). Since the abstract basis need not be
  $\langle\,,\rangle$-orthonormal, $B$ is a \emph{general} constant symmetric
  positive-definite matrix, not $\delta$: do Carmo's $\delta_{ij}$ of $(1)$ is
  replaced by $B_{ij}$ in the abstract chart frame, and the whole §3 conformal
  computation goes through with $\delta\rightsquigarrow B$, the extra factors
  collapsing by the inverse-Gram identity $\sum_{ij}B^{ij}c_ic_j=1$ (for
  $c_i=\langle e_n,\text{finBasis}_i\rangle$).
\end{lemma}
\begin{proof}
  \leanok
  The trivialisation of the tangent bundle of an open subset of $\mathbb R^n$ is
  the identity on fibres, so $\partial_i(x)=$ finBasis$_i$ and
  $G_{ij}(x)=\langle\partial_i,\partial_j\rangle_g=x_n^{-2}\langle\text{finBasis}_i,\text{finBasis}_j\rangle=x_n^{-2}B_{ij}$.
  The Gram matrix at the basepoint is positive definite, so $\det B$ is a unit and
  the inverse of the scalar multiple $x_n^{-2}B$ is $x_n^{2}B^{-1}$.
\end{proof}

\begin{lemma}[hyperbolic chart Christoffel symbols and their derivatives, closed form]
  \label{lem:dc-ch8-3-hyperbolic-christoffel}
  \lean{Riemannian.Hyperbolic.hyperbolic_chartChristoffel,
    Riemannian.Hyperbolic.hyperbolic_partialDeriv_chartChristoffel,
    Riemannian.Hyperbolic.hyperbolic_partialDeriv_chartGramOnE}
  \leanok
  \uses{lem:dc-ch8-3-hyperbolic-chart-gram}
  In the abstract chart frame the Christoffel symbols of the hyperbolic metric are,
  in closed form,
  $$\Gamma^k_{ij}(y)=-f_i\delta^k_j-f_j\delta^k_i+d^k B_{ij},\qquad
    f_i=\frac{c_i}{y_n},\ c_i=(\text{finBasis}_i)_n,\ d^k=\frac{\sum_\ell B^{k\ell}c_\ell}{y_n},$$
  the constant abstract-basis Gram $B$ (\cref{lem:dc-ch8-3-hyperbolic-chart-gram})
  playing the role of do Carmo's $\delta$
  ($\Gamma^k_{ij}=-\delta_{jk}f_i-\delta_{ki}f_j+\delta_{ij}f_k$ with $\delta\rightsquigarrow B$).
  Because every term carries the single factor $1/y_n$, one may write
  $\Gamma^k_{ij}(y)=K^k_{ij}\,(y_n)^{-1}$ with $K^k_{ij}$ the ($y$-independent)
  numerator, so the abstract-chart-frame partial derivatives are
  $$\partial_m\Gamma^k_{ij}(y)=K^k_{ij}\cdot\bigl(-c_m\,(y_n)^{-2}\bigr).$$
  This is exactly the $\Gamma$ and $\partial\Gamma$ data consumed by the chart-frame
  curvature expansion \cref{lem:dc-ch8-3-chart-curvature}.
\end{lemma}
\begin{proof}
  \leanok
  The chart Gram is $G_{ij}(y)=(y_n^2)^{-1}B_{ij}$ and its inverse
  $G^{k\ell}(y)=y_n^2 B^{k\ell}$ (\cref{lem:dc-ch8-3-hyperbolic-chart-gram}). The
  distinguished coordinate $y_n=y_e$ has derivative
  $\partial_m y_n=(\text{finBasis}_m)_n=c_m$ along $\text{finBasis}_m$, so
  $\partial_m\bigl((y_n^2)^{-1}\bigr)=-2\,y_n^{-3}c_m$ and hence
  $\partial_m G_{ij}=B_{ij}\cdot(-2\,y_n^{-3}c_m)$. Substituting into the Koszul
  formula
  $\Gamma^k_{ij}=\tfrac12\sum_\ell G^{k\ell}(\partial_iG_{j\ell}+\partial_jG_{i\ell}-\partial_\ell G_{ij})$
  and contracting $\sum_\ell B^{k\ell}B_{\ell j}=\delta^k_j$,
  $\sum_\ell B^{k\ell}B_{\ell i}=\delta^k_i$ (from $B^{-1}B=1$) collapses the factors
  $y_n^2\cdot y_n^{-3}=y_n^{-1}$, giving the closed form. Since only the $1/y_n$ factor
  depends on $y$, $\partial_m\Gamma^k_{ij}=K^k_{ij}\,\partial_m\bigl(y_n^{-1}\bigr)
  =K^k_{ij}\,(-c_m y_n^{-2})$.
\end{proof}

\begin{lemma}[inverse-Gram/coordinate identity $\sum_{ij}B^{ij}c_ic_j=1$]
  \label{lem:dc-ch8-3-inverse-gram-coord}
  \lean{Riemannian.Hyperbolic.finBasis_inv_gram_coord}
  \leanok
  \uses{lem:dc-ch8-3-hyperbolic-chart-gram}
  For the constant abstract-basis Gram matrix
  $B_{ij}=\langle\text{finBasis}_i,\text{finBasis}_j\rangle$ with inverse
  $B^{ij}=(B^{-1})_{ij}$, and $c_i=(\text{finBasis}_i)_n$ the $n$-th coordinate of
  the $i$-th basis vector,
  $$\sum_{i,j} B^{ij}\,c_i\,c_j = 1.$$
  Equivalently $c_i=\langle\text{finBasis}_i, w_n\rangle$ for the standard unit
  vector $w_n=\text{basisFun}_n$, and the $B^{-1}$-contraction reconstructs
  $\langle w_n, w_n\rangle=1$. This is exactly the term
  $|\mathrm d\sigma|^2_{g_0}=1$ that makes the §3 conformal cancellation
  $T=\tfrac12\varphi g_0$ (hence $K\equiv-1$) go through with the \emph{general}
  non-orthonormal abstract Gram $B$ in place of $\delta$.
\end{lemma}
\begin{proof}
  \leanok
  Write $w_n=\sum_i a_i\,\text{finBasis}_i$ in the abstract basis (the
  \texttt{repr} coordinates). Then $c_i=\langle\text{finBasis}_i,w_n\rangle=\sum_j
  B_{ij}a_j$, so $\sum_j B^{ij}c_j=(B^{-1}B a)_i=a_i$ and
  $\sum_{ij}B^{ij}c_ic_j=\sum_i a_i c_i=\langle w_n,w_n\rangle=1$.
\end{proof}

\begin{lemma}[the pointwise $(0,4)$ curvature form and its chart-frame value]
  \label{lem:dc-ch8-3-pointwise-curvature-form}
  \lean{Riemannian.AffineConnection.curvatureFormAt,
    Riemannian.AffineConnection.isAlgCurvatureForm_curvatureFormAt,
    Riemannian.leviCivita_curvatureFormAt_chartFrame}
  \leanok
  \uses{lem:dc-ch8-3-chart-curvature}
  The curvature operator $R(X,Y)Z$ is a tensor: its value at $p$ depends only on
  $X(p),Y(p),Z(p)$ (do Carmo Ch.~4, Prop.~2.2 read pointwise). Lowering by the
  metric gives the pointwise $(0,4)$ form $\langle R(x,y)z,t\rangle_g$ on
  $T_pM$, which for the Levi-Civita connection is an \emph{algebraic curvature
  form} (multilinear, antisymmetric in each pair, first Bianchi — do Carmo
  Prop.~2.5). In the chart frame $\partial_a=\text{finBasis}_a$ its value is the
  coordinate curvature coefficient contracted with the chart Gram,
  $$\langle R(\partial_i,\partial_j)\partial_k,\partial_\ell\rangle_g
    = \sum_m R^m_{ijk}\,G_{m\ell},\qquad
    R^m_{ijk}=\partial_j\Gamma^m_{ik}-\partial_i\Gamma^m_{jk}
      +\sum_s(\Gamma^s_{ik}\Gamma^m_{js}-\Gamma^s_{jk}\Gamma^m_{is}),$$
  the metric-lowered avatar of \cref{lem:dc-ch8-3-chart-curvature}. This bridges
  the \emph{intrinsic} manifold curvature to the coordinate Christoffel data, and
  is the input to the all-planes identification \cref{cor:dc-ch4-3-5}: it makes the
  chart coefficients literally the components of an algebraic curvature form on
  $T_pM$, ready for \texttt{ext\_basis}/\texttt{eq\_kronecker\_iff\_const}.
\end{lemma}
\begin{proof}
  \leanok
  Tensoriality is the pointwise reading of the $\mathcal D(M)$-multilinearity of
  $R$: a field vanishing at $p$ decomposes as $\sum_i f_i W_i+\tau$ with
  $f_i(p)=0$ and $\tau$ vanishing near $p$, and homogeneity of $R$ in that slot
  kills both parts, so $R(\cdot)$ against a field vanishing at $p$ vanishes at
  $p$; hence $R(X,Y)Z|_p$ depends only on the values at $p$. The
  algebraic-curvature-form axioms transport slotwise from the field-level
  symmetries (Prop.~2.5, via the Levi-Civita property). The chart-frame value
  pairs the frame expansion of $R(\partial_i,\partial_j)\partial_k$
  (\cref{lem:dc-ch8-3-chart-curvature}) against $\partial_\ell$, distributing the
  metric over the finite sum.
\end{proof}

\begin{proposition}[$H^n$ has constant sectional curvature $-1$]
  \label{prop:dc-ch8-3-const-curv}
  \lean{Riemannian.Hyperbolic.hyperbolic_isConstantCurvature,
    Riemannian.Hyperbolic.hyperbolicMetric_sectionalCurvature_eq_neg_one,
    Riemannian.Hyperbolic.hyperbolic_curvatureFormAt_eq,
    Riemannian.Hyperbolic.hyperbolic_curvatureFormAt_frame}
  \leanok
  \uses{lem:dc-ch8-3-conformal-curvature,lem:dc-ch8-3-chart-curvature,lem:dc-ch8-3-pointwise-curvature-form,lem:dc-ch8-3-hyperbolic-chart-gram,lem:dc-ch8-3-hyperbolic-christoffel,lem:dc-ch8-3-inverse-gram-coord,cor:dc-ch4-3-5,def:dc-ch8-3-hyperbolic-space}
  The hyperbolic space $H^n$, with the metric $g_{ij}=\delta_{ij}/x_n^2$, has
  constant sectional curvature equal to $-1$: for every point and every
  $2$-plane $\sigma\subset T_pH^n$, $K(p,\sigma)=-1$.
\end{proposition}
\begin{proof}
  \leanok
  Work in the abstract chart frame $\partial_a=\text{finBasis}_a$ of $T_pH^n$. By
  \cref{lem:dc-ch8-3-pointwise-curvature-form} the pointwise $(0,4)$ curvature form
  of the Levi-Civita connection is, in this frame,
  $\langle R(\partial_i,\partial_j)\partial_k,\partial_\ell\rangle_g=\sum_m
  R^m_{ijk}G_{m\ell}$, with $R^m_{ijk}=\partial_j\Gamma^m_{ik}-\partial_i\Gamma^m_{jk}
  +\sum_s(\Gamma^s_{ik}\Gamma^m_{js}-\Gamma^s_{jk}\Gamma^m_{is})$ and $G_{m\ell}$
  the chart Gram matrix. Insert the closed forms
  \cref{lem:dc-ch8-3-hyperbolic-christoffel}, $\Gamma^k_{ij}(y)=y_n^{-1}K^k_{ij}$
  and $\partial_m\Gamma^k_{ij}(y)=-c_m\,y_n^{-2}K^k_{ij}$ with the \emph{constant}
  coefficient $K^k_{ij}=-c_i\delta^k_j-c_j\delta^k_i+D^kB_{ij}$
  ($c_i=(\text{finBasis}_i)_n$, $D^k=\sum_\ell B^{k\ell}c_\ell$), and the chart Gram
  \cref{lem:dc-ch8-3-hyperbolic-chart-gram}, $G_{m\ell}=y_n^{-2}B_{m\ell}$ with $B$
  the constant abstract-basis Gram matrix. Every term carries a factor $y_n^{-4}$,
  which cancels against the right-hand side, and the identity reduces to the purely
  finite contraction $\sum_m R^m_{ijk}B_{m\ell}=-(B_{ik}B_{j\ell}-B_{jk}B_{i\ell})$
  in the constants $c,B,D$. This contraction is do Carmo's conformal cancellation
  $T=\tfrac12\varphi g_0$: expanding the Christoffel products and contracting each
  index sum with the ``lowered'' identity $\sum_m K^m_{ab}B_{m\ell}=-c_aB_{b\ell}
  -c_bB_{a\ell}+c_\ell B_{ab}$ (which uses $\sum_m D^mB_{m\ell}=c_\ell$), every
  monomial cancels except $-(B_{ik}B_{j\ell}-B_{jk}B_{i\ell})$, the last step using
  the single nontrivial input $\sum_{ij}B^{ij}c_ic_j=1$
  (\cref{lem:dc-ch8-3-inverse-gram-coord}). Hence
  $\langle R(\partial_i,\partial_j)\partial_k,\partial_\ell\rangle_g
  =-(G_{ik}G_{j\ell}-G_{jk}G_{i\ell})$ on the chart frame. Both the pointwise
  curvature form and $-1$ times the standard curvature form
  $\langle x,z\rangle\langle y,t\rangle-\langle y,z\rangle\langle x,t\rangle$ are
  algebraic curvature forms (\cref{lem:dc-ch8-3-pointwise-curvature-form}) that
  agree on this basis frame, so by the basis-determination step of
  \cref{cor:dc-ch4-3-5} (\texttt{ext\_basis}) they agree on \emph{all} of $T_pH^n$:
  $\langle R(x,y)z,t\rangle_g=-(\langle x,z\rangle_g\langle y,t\rangle_g
  -\langle y,z\rangle_g\langle x,t\rangle_g)$. This is constant curvature $-1$ at the
  field level; dividing the diagonal $\langle R(x,y)x,y\rangle_g=-|x\wedge y|^2$ by
  $|x\wedge y|^2$ gives $K(p,\sigma)=-1$ for every $2$-plane
  $\sigma=\mathrm{span}\{x,y\}$.
\end{proof}

\begin{lemma}[Christoffel contraction of the hyperbolic metric, closed form]
  \label{lem:dc-ch8-3-christoffel-contraction}
  \lean{Riemannian.Hyperbolic.hyperbolic_chartChristoffelContraction_eq}
  \leanok
  \uses{lem:dc-ch8-3-hyperbolic-christoffel,lem:dc-ch8-3-inverse-gram-coord}
  Contracting the closed-form chart Christoffel symbols
  $\Gamma^k_{ij}=-\delta_{jk}f_i-\delta_{ki}f_j+\delta_{ij}f_k$
  ($f_i=(\partial_i)_n/x_n$) of \cref{lem:dc-ch8-3-hyperbolic-christoffel} against
  two tangent vectors $v,w\in T_pH^n$ collapses, via the inverse-Gram/coordinate
  identity \cref{lem:dc-ch8-3-inverse-gram-coord}, to the basis-independent
  Euclidean-valued expression
  $$
  \Gamma(v,w)(y)=\nabla_v w
    =\frac{1}{y_n}\Big(\langle v,w\rangle\, e_n-v_n\,w-w_n\,v\Big),
  $$
  where $y_n$ is the distinguished coordinate of the base point, $v_n=\langle v,e_n\rangle$,
  $e_n$ is the distinguished unit vector, and $\langle\,,\rangle$ is the ambient
  Euclidean inner product. This is do Carmo's covariant derivative of $H^n$ read in
  the trivial chart of the open half-space.
\end{lemma}
\begin{proof}
  \leanok
  In the abstract chart frame the contraction is
  $\Gamma(v,w)(y)=\sum_k\big(\sum_{ij}\Gamma^k_{ij}(y)\,v^i w^j\big)\partial_k$ with
  $v^i,w^j$ the frame coordinates. Substituting the closed form
  \cref{lem:dc-ch8-3-hyperbolic-christoffel} and summing the two Kronecker terms
  gives $-\big(\sum_i (\partial_i)_n v^i\big)w-\big(\sum_j(\partial_j)_n w^j\big)v
  =-v_n w-w_n v$ (using $\sum_i(\partial_i)_n v^i=v_n$), while the $\delta_{ij}f_k$
  term contracts to $\langle v,w\rangle\sum_k(\sum_\ell B^{k\ell}c_\ell)\partial_k
  =\langle v,w\rangle\,e_n$ by \cref{lem:dc-ch8-3-inverse-gram-coord}; the common
  factor $1/y_n$ is carried throughout.
\end{proof}

\begin{lemma}[geodesic equation of $H^n$ in Euclidean form]
  \label{lem:dc-ch8-3-geodesic-equation}
  \lean{Riemannian.Hyperbolic.hyperbolic_isGeodesic_of}
  \leanok
  \uses{lem:dc-ch8-3-christoffel-contraction}
  A twice-differentiable curve $u:\mathbb R\to\mathbb R^n$ staying in the
  half-space ($u_n(t)>0$) whose acceleration satisfies the ordinary differential
  equation
  $$
  u''(t)+\frac{1}{u_n(t)}\Big(\langle u'(t),u'(t)\rangle\, e_n-2\,u'_n(t)\,u'(t)\Big)=0
  $$
  is a geodesic of $H^n$. This is the geodesic equation $u''+\Gamma(u',u')=0$ read
  through \cref{lem:dc-ch8-3-christoffel-contraction}; because the chart of the open
  half-space is the inclusion, the moving-foot chart curve of the intrinsic geodesic
  predicate is literally $u$.
\end{lemma}
\begin{proof}
  \leanok
  For an open subset of $\mathbb R^n$ the extended chart at any base point is the
  inclusion, so the chart-local curve $s\mapsto\varphi_{\gamma t}(\gamma s)$ equals
  $u$ for every $t$; its first and second derivatives are $u'$ and $u''$. The
  intrinsic moving-foot geodesic condition at time $t$ is
  $u''(t)+\Gamma(u'(t),u'(t))(u(t))=0$, and substituting the closed form
  \cref{lem:dc-ch8-3-christoffel-contraction} (with $v=w=u'(t)$, so the two mixed
  terms coincide) gives exactly the displayed equation.
\end{proof}

\begin{lemma}[vertical lines are geodesics of $H^n$]
  \label{lem:dc-ch8-3-vertical-geodesic}
  \lean{Riemannian.Hyperbolic.hyperbolic_vertical_isGeodesic}
  \leanok
  \uses{lem:dc-ch8-3-geodesic-equation}
  For a horizontal base point $c$ (with $c_n=0$) and a constant $a>0$, the affinely
  parametrised vertical line
  $$
  t\longmapsto c+(a\,e^{t})\,e_n,\qquad t\in\mathbb R,
  $$
  perpendicular to the boundary hyperplane $\{x_n=0\}$, is a geodesic of $H^n$. Its
  distinguished coordinate $u_n(t)=a\,e^{t}$ solves do Carmo's $u_n''\,u_n=(u_n')^2$.
\end{lemma}
\begin{proof}
  \leanok
  Here $u(t)=c+(a e^{t})e_n$, $u'(t)=u''(t)=(a e^{t})e_n$, $u_n(t)=u'_n(t)=a e^{t}$
  and $\langle u',u'\rangle=(a e^{t})^2$. The geodesic equation
  \cref{lem:dc-ch8-3-geodesic-equation} reads
  $(a e^{t})e_n+(a e^{t})^{-1}\big((a e^{t})^2 e_n-2(a e^{t})^2 e_n\big)
  =(a e^{t})e_n-(a e^{t})e_n=0$, which holds identically.
\end{proof}

\begin{lemma}[semicircles perpendicular to $\partial H^n$ are geodesics of $H^n$]
  \label{lem:dc-ch8-3-semicircle-geodesic}
  \lean{Riemannian.Hyperbolic.hyperbolic_semicircle_isGeodesic}
  \leanok
  \uses{lem:dc-ch8-3-geodesic-equation}
  For a center $m$ on the boundary hyperplane ($m_n=0$), a unit horizontal vector
  $\hat u$ ($\hat u_n=0$, $\langle\hat u,\hat u\rangle=1$) and a radius $r>0$, the
  affinely parametrised semicircle
  $$
  t\longmapsto m+(r\tanh t)\,\hat u+(r\,\mathrm{sech}\,t)\,e_n,\qquad t\in\mathbb R,
  $$
  (with $\tanh t=\sinh t/\cosh t$, $\mathrm{sech}\,t=1/\cosh t$), perpendicular to
  the boundary $\{x_n=0\}$ and centered on it, is a geodesic of $H^n$.
\end{lemma}
\begin{proof}
  \leanok
  With $u(t)=m+(r\tanh t)\hat u+(r\,\mathrm{sech}\,t)e_n$ one has
  $u'(t)=(r\,\mathrm{sech}^2 t)\hat u-(r\tanh t\,\mathrm{sech}\,t)e_n$; using
  $\mathrm{sech}'=-\tanh\,\mathrm{sech}$, $\tanh'=\mathrm{sech}^2$, and
  $\cosh^2-\sinh^2=1$ one computes $u_n=r\,\mathrm{sech}\,t$,
  $\langle u',u'\rangle=r^2\mathrm{sech}^2 t$, and both the $\hat u$- and
  $e_n$-components of $u''+u_n^{-1}(\langle u',u'\rangle e_n-2u'_n u')$ vanish
  identically. Hence \cref{lem:dc-ch8-3-geodesic-equation} applies.
\end{proof}

\begin{proposition}[geodesics of $H^n$]
  \label{prop:dc-ch8-3-1}
  \dcref{ch8:3.1}
  \lean{Riemannian.Hyperbolic.hyperbolic_geodesic_classification}
  \leanok
  \uses{ex:dc-ch3-3-10,lem:dc-ch8-3-geodesic-equation,lem:dc-ch8-3-vertical-geodesic,
    lem:dc-ch8-3-semicircle-geodesic}
  The straight lines perpendicular to the hyperplane $x_n=0$, and the circles of
  $H^n$ whose planes are perpendicular to the hyperplane $x_n=0$ and whose centers
  are in this hyperplane, are the geodesics of $H^n$.
\end{proposition}
\begin{proof}
  \leanok
  That each stated curve is a geodesic is the forward direction, already established:
  the vertical lines by \cref{lem:dc-ch8-3-vertical-geodesic} and the semicircles
  perpendicular to $\{x_n=0\}$ with centre on it by
  \cref{lem:dc-ch8-3-semicircle-geodesic}, both through the geodesic equation
  \cref{lem:dc-ch8-3-geodesic-equation}. It remains to prove the converse inclusion:
  that these are \emph{all} the geodesics.

  Let $\gamma$ be a geodesic of $H^n$, and read it through the inclusion chart of the
  open half-space, so that $\gamma$ is an ambient curve with $\gamma(0)=p$ and chart
  velocity $\gamma'(0)=v\in\mathbb R^n$. Decompose $v=v_h+v_n\,e_n$, where
  $e_n=\partial/\partial x_n$, $v_n=\langle v,e_n\rangle$, and $v_h$ is the horizontal
  part (so $\langle v_h,e_n\rangle=0$). Through $(p,v)$ passes exactly one member of
  the stated family:
  \begin{itemize}
    \item if $v_h=0$, the reparametrised vertical line
      $s\mapsto c+p_n\,e^{Bs}\,e_n$, with $c=p-p_n e_n$ (so $c_n=0$), $p_n=\langle
      p,e_n\rangle>0$ and $B=v_n/p_n$;
    \item if $v_h\neq0$, the semicircle
      $s\mapsto m+r\tanh(\sigma s+s_0)\,\hat u+r\,\operatorname{sech}(\sigma s+s_0)\,e_n$
      in the plane $\{e_n,\hat u\}$, with $\hat u=v_h/|v_h|$, and centre, radius,
      speed and phase determined by $(p,v)$: writing $s_0=\operatorname{arsinh}(-v_n/|v_h|)$,
      $r=p_n\cosh s_0$, $\sigma=|v_h|\cosh s_0/p_n$ and
      $m=(p-p_n e_n)+\bigl(p_n v_n/|v_h|\bigr)\hat u$, so that $m_n=0$, $\hat u_n=0$,
      $|\hat u|=1$, $r>0$.
  \end{itemize}
  A direct computation (using $\sinh s_0=-v_n/|v_h|$ and
  $\cosh^2-\sinh^2=1$) shows that in either case this family member is a geodesic
  passing through $p$ with chart velocity $v$ at $s=0$. By the uniqueness of geodesics
  through a given point with a given velocity (do Carmo Ch.~3, Theorem 2.2, intrinsic
  form), $\gamma$ coincides with it, so $\gamma$ traces a vertical line or such a
  semicircle. This is the required converse.

  Since every geodesic is thus contained in a plane perpendicular to $x_n=0$, and such
  planes are isometric to the hyperbolic plane, the completeness of $H^n$ follows from
  that of the hyperbolic plane (cf. Exercise 10, Chap.~7). Another model of hyperbolic
  space is given in Exercise 3 of this chapter. The book's alternative route reduces to
  the $x_1x_n$-plane through \cref{ex:dc-ch3-3-10} via an isometry of $\mathbb R^n$ in
  the variables $x_1,\dots,x_{n-1}$.
\end{proof}

\section{Space forms}

Now we are able to prove the main theorem of this chapter. As always, $M^n$
denotes a manifold of dimension $n$.

\begin{theorem}[universal covering of a space of constant curvature]
  \label{thm:dc-ch8-4-1}
  \dcref{ch8:4.1}
  \notready
  Let $M^n$ be a complete Riemannian manifold with constant sectional curvature $K$.
  Then the universal covering $\tilde{M}$ of $M$, with the covering metric, is
  isometric to:
  \begin{enumerate}
    \item[(a)] $H^n$, if $K=-1$;
    \item[(b)] $\mathbb{R}^n$, if $K=0$;
    \item[(c)] $S^n$, if $K=1$.
  \end{enumerate}
\end{theorem}
\begin{proof}
  \uses{thm:dc-ch8-2-1,lem:dc-ch8-4-2,prop:dc-ch8-3-const-curv,
    lem:dc-ch7-3-1-covering-diffeo,lem:dc-ch7-3-1-covering-simplyconnected}
  $\tilde{M}$ is a simply connected, complete Riemannian
  manifold with constant sectional curvature $K$. Consider first the cases (a) and
  (b) and denote, for convenience, both $H^n$ and $\mathbb{R}^n$ by $\triangle$. Fix
  points $p\in\triangle$, $\tilde{p}\in\tilde{M}$ and a linear isometry
  $i:T_p(\triangle)\to T_{\tilde{p}}(\tilde{M})$. Consider the map
  $f=\exp_{\tilde{p}}\circ\,i\circ\exp_p^{-1}:\triangle\to\tilde{M}$. Since $\triangle$
  and $\tilde{M}$ are complete with non-positive sectional curvature, $f$ is
  well-defined. From the Theorem of Cartan, $f$ is a local isometry. From Lemma 3.3
  of Chap. 7, $f$ is a diffeomorphism, which proves (a) and (b).

  For case (c), fix $p\in S^n$, $\tilde{p}\in\tilde{M}$ and a linear isometry
  $i:T_p(S^n)\to T_{\tilde{p}}(\tilde{M})$. Let $q\in S^n$ be the antipodal point of
  $p$ and define $f=\exp_{\tilde{p}}\circ\,i\circ\exp_p^{-1}:S^n-\{q\}\to\tilde{M}$.
  From the Theorem of Cartan, $f$ is a local isometry. Choose $p'\in S^n$,
  $p'\neq p$, $p'\neq q$, set $\tilde{p}'=f(p')$, $i'=df_{p'}$ and define
  $f'=\exp_{\tilde{p}'}\circ\,i'\circ\exp_{p'}^{-1}:S^n-\{q'\}\to\tilde{M}$, where
  $q'$ is the antipodal point of $p'$. On the connected set
  $W=S^n-(\{q\}\cup\{q'\})$ we have $p'\in W$ and $f(p')=\tilde{p}'=f'(p')$,
  $df_{p'}=i'=df'_{p'}$, so by Lemma 4.2, $f=f'$ on $W$. We can then define
  $g:S^n\to\tilde{M}$ by $g(r)=f(r)$ if $r\in S^n-\{q\}$ and $g(r)=f'(r)$ if
  $r\in S^n-\{q'\}$. Then $g$ is a local isometry, hence a local diffeomorphism. By
  the compactness of $S^n$, $g$ is a covering map, and since $\tilde{M}$ is simply
  connected, $g$ is a diffeomorphism. Therefore $g$ is an isometry. $\square$
\end{proof}

\begin{lemma}[a local isometry pushes geodesics forward]
  \label{lem:dc-ch8-4-2-geodesic-push}
  \lean{Riemannian.Geodesic.isGeodesicOn_comp_of_isLocalDiffeomorph}
  \leanok
  \uses{def:dc-ch1-2-3,def:dc-ch3-2-1}
  Let $f:M\to N$ be a local isometry and let $\gamma$ be a geodesic of $M$ defined on
  an open set $s$ of times. Then $f\circ\gamma$ is a geodesic of $N$ on $s$.
\end{lemma}
\begin{proof}
  \leanok
  Being a local isometry, $f$ identifies the metric of $M$ with the pullback $f^*g_N$,
  so $\gamma$ satisfies the $f^*g_N$-geodesic equation. The geodesic equation is a
  pointwise condition, and under $f$ the geodesic operator transforms by the
  differential $df$ alone: from $df\big(\gamma''+\Gamma^{f^*g_N}(\gamma',\gamma')\big)
  =(f\circ\gamma)''+\Gamma^{g_N}\big((f\circ\gamma)',(f\circ\gamma)'\big)$ and the
  vanishing of the left-hand argument one reads off that $f\circ\gamma$ satisfies the
  $g_N$-geodesic equation at each time of $s$. No injectivity of $f$ is needed.
\end{proof}

\begin{lemma}[the agreement locus is open and closed]
  \label{lem:dc-ch8-4-2-agree-locus}
  \lean{Riemannian.AgreeLocus,
    Riemannian.mem_agreeLocus_iff,
    Riemannian.isClosed_agreeLocus,
    Riemannian.eq_of_pointDatum_of_preconnected}
  \leanok
  Let $f_1,f_2:M\to N$ be differentiable, $M$ connected, and let
  $$
  A=\{x\in M;\ f_1(x)=f_2(x)\ \text{and}\ (df_1)_x=(df_2)_x\}
  $$
  be their \emph{agreement locus}. Then $A$ is closed. If moreover $A$ is open — for
  instance if agreement to first order at a point always propagates to a neighborhood
  of that point — and $A\neq\emptyset$, then $f_1=f_2$.
\end{lemma}
\begin{proof}
  \leanok
  A point $x$ lies in $A$ exactly when the tangent maps $df_1$ and $df_2$ agree over
  $x$, since the tangent map of $f$ sends $(x,v)$ to $(f(x),(df)_x v)$: the base
  components give $f_1(x)=f_2(x)$ and the fibre components $(df_1)_x=(df_2)_x$. So $A$
  is the complement of the image, under the bundle projection $TM\to M$, of the locus
  where the two tangent maps differ. The tangent maps are continuous and $TN$ is
  Hausdorff, so that locus is open; the bundle projection is an open map, so its image
  is open and $A$ is closed. If $A$ is also open and nonempty then, $M$ being
  connected, $A=M$, i.e. $f_1=f_2$.
\end{proof}

\begin{lemma}[first-order agreement propagates]
  \label{lem:dc-ch8-4-2-propagation}
  \lean{Riemannian.eventuallyEq_of_pointDatum}
  \leanok
  \uses{lem:dc-ch8-4-2-geodesic-push,prop:dc-ch3-2-9,def:dc-ch3-expmap,
    lem:dc-ch3-expmap-ray,thm:dc-ch3-2-2}
  Let $f_1,f_2:M\to N$ be local isometries and let $q\in M$ satisfy $f_1(q)=f_2(q)$ and
  $(df_1)_q=(df_2)_q$. Then $f_1=f_2$ on a neighborhood of $q$.
\end{lemma}
\begin{proof}
  \leanok
  \uses{lem:dc-ch8-4-2-geodesic-push,prop:dc-ch3-2-9,def:dc-ch3-expmap,
    lem:dc-ch3-expmap-ray,thm:dc-ch3-2-2}
  By \cref{prop:dc-ch3-2-9}, $\exp_q$ carries a ball $B_\varepsilon(0)\subset T_qM$
  diffeomorphically onto an open neighborhood of $q$; so it suffices to prove
  $f_1(\exp_q u)=f_2(\exp_q u)$ for all sufficiently small $u$. Fix such a $u$ and let
  $\gamma_u(t)=\exp_q(tu)$, a geodesic of $M$ on an interval $(-b,b)$ with $b>1$, with
  $\gamma_u(0)=q$ and $\gamma_u'(0)=u$ (\cref{lem:dc-ch3-expmap-ray}). By
  \cref{lem:dc-ch8-4-2-geodesic-push} both $f_1\circ\gamma_u$ and $f_2\circ\gamma_u$ are
  geodesics of $N$ on $(-b,b)$. At $t=0$ they have the same initial point, since
  $f_1(q)=f_2(q)$, and the same initial velocity, since
  $(df_1)_q u=(df_2)_q u$. By uniqueness of geodesics with given initial conditions
  (\cref{thm:dc-ch3-2-2}) they coincide on the connected interval $(-b,b)$; evaluating
  at $t=1$ gives $f_1(\exp_q u)=f_2(\exp_q u)$.
\end{proof}

\begin{lemma}[uniqueness of a local isometry from a point datum]
  \label{lem:dc-ch8-4-2}
  \dcref{ch8:4.2}
  \lean{Riemannian.eq_of_pointDatum_of_localIsometry}
  \leanok
  \uses{def:dc-ch1-2-3}
  Let $f_i:M\to N$, $i=1,2$, be two local isometries of the (connected) Riemannian
  manifold $M$ to the Riemannian manifold $N$, both of dimension $n$ and modelled on a
  common Euclidean space. Suppose that there exists a point $p\in M$ such that
  $f_1(p)=f_2(p)$ and $(df_1)_p=(df_2)_p$. Then $f_1=f_2$.
\end{lemma}
\begin{proof}
  \leanok
  \uses{lem:dc-ch8-4-2-agree-locus,lem:dc-ch8-4-2-propagation}
  Consider the agreement locus
  $A=\{x\in M;\ f_1(x)=f_2(x)\ \text{and}\ (df_1)_x=(df_2)_x\}$. It contains $p$, so it
  is nonempty; it is closed by \cref{lem:dc-ch8-4-2-agree-locus}; and it is open,
  because by \cref{lem:dc-ch8-4-2-propagation} each of its points has a neighborhood on
  which $f_1$ and $f_2$ agree, and on such a neighborhood the differentials agree too.
  Since $M$ is connected, $A=M$, that is, $f_1=f_2$.
\end{proof}

The complete manifolds with constant sectional curvature are called \emph{space forms}.
The last theorem reduces the determination of all the space forms to a problem in
group theory. Recall: a group $G$ \emph{acts} on a set $M$ if there is a map
$G\times M\to M$, $(g,x)\mapsto gx$, with $ex=x$ and $(g_1 g_2)x=g_1(g_2 x)$; $G$
acts \emph{freely} if $gx=x$ implies $g=e$; the \emph{orbit} of $x$ is $Gx=\{gx;\ g\in G\}$.
The action is \emph{transitive} if $Gx=M$. If $M$ is a topological space, $G$ (a group
of homeomorphisms) acts in a \emph{totally discontinuous manner} if every $x\in M$ has a
neighborhood $U$ such that $g(U)\cap U=\emptyset$ for all $g\in G$, $g\neq e$; in
this case $\pi:M\to M/G$ is a regular covering map and $G$ is the group of covering
transformations.

If $M$ is a Riemannian manifold and $\Gamma$ is a subgroup of the group of
isometries acting in a totally discontinuous manner, then $M/\Gamma$ has a
differentiable structure in which $\pi:M\to M/\Gamma$ is a local diffeomorphism,
and a Riemannian metric making $\pi$ a local isometry. Given $p\in M/\Gamma$,
choose $\tilde p\in\pi^{-1}(p)$; for $u,v\in T_p(M/\Gamma)$, define

$$
\langle u,v\rangle=\langle d\pi^{-1}(u),d\pi^{-1}(v)\rangle_{\tilde{p}},
$$

Since $\pi$ is a regular covering, $\Gamma$ is transitive on $\pi^{-1}(p)$, and
because $\Gamma$ consists of isometries the definition is independent of the
choice of $\tilde p$. This is called the \emph{metric on $M/\Gamma$ induced by the
covering $\pi$}. $M/\Gamma$ is complete iff $M$ is, and has constant curvature iff
$M$ does. Taking $M=S^n$,
$\mathbb{R}^n$ or $H^n$, we conclude that $M/\Gamma$ is a complete manifold of
constant curvature $1$, $0$ or $-1$, respectively.

\begin{proposition}[every space form is a quotient of a model]
  \label{prop:dc-ch8-4-3}
  \dcref{ch8:4.3}
  \notready
  Let $M$ be a complete Riemannian manifold with constant sectional curvature $K$
  ($1,0,-1$). Then $M$ is isometric to $\tilde{M}/\Gamma$, where $\tilde{M}$ is $S^n$
  (if $K=1$), $\mathbb{R}^n$ (if $K=0$) or $H^n$ (if $K=-1$), $\Gamma$ is a subgroup
  of the group of isometries of $\tilde{M}$ which acts in a totally discontinuous
  manner on $\tilde{M}$, and the metric on $\tilde{M}/\Gamma$ is induced from the
  covering $\pi:\tilde{M}\to\tilde{M}/\Gamma$.
\end{proposition}
\begin{proof}
  \uses{thm:dc-ch8-4-1,lem:dc-ch7-3-1-covering-simplyconnected}
  Consider the universal covering $p:\tilde{M}\to M$ with the covering
  metric. Let $\Gamma$ be the group of covering transformations of $p$. Then
  $\Gamma$ is a subgroup of the isometries of $\tilde{M}$ and acts in a totally
  discontinuous manner. Introduce on $\tilde{M}/\Gamma$ the metric induced by
  $\pi:\tilde{M}\to\tilde{M}/\Gamma$. Since the covering $p$ is regular, given
  $\tilde{x},\tilde{y}\in\tilde{M}$, $p(\tilde{x})=p(\tilde{y})$ iff
  $\Gamma\tilde{x}=\Gamma\tilde{y}$, equivalently $\pi(\tilde{x})=\pi(\tilde{y})$.
  The equivalence classes given by $p$ and $\pi$ are the same, which induces a
  bijection $\xi:M\to\tilde{M}/\Gamma$ such that $\pi=\xi\circ p$. Since $\pi$ and
  $p$ are local isometries, $\xi$ is a local isometry, and being a bijection, is an
  isometry of $M$ onto $\tilde{M}/\Gamma$.
\end{proof}

The last proposition reduces the problem of finding all the space forms to that of
determining all the subgroups of the group of isometries that act in a totally
discontinuous manner on each of the simply connected models. The spherical problem
($\tilde{M}=S^n$) was solved during the sixties (Wolf). Here we mention only two
interesting facts.

\begin{proposition}[even-dimensional space forms of $K=1$]
  \label{prop:dc-ch8-4-4}
  \dcref{ch8:4.4}
  \notready
  Let $M^n$ be a complete Riemannian manifold of even dimension $n=2m$, with constant
  sectional curvature $K=1$. Then $M^n$ is isometric to the sphere $S^n$ or the real
  projective space $P^n$ of the same dimension.
\end{proposition}
\begin{proof}
  \uses{prop:dc-ch8-4-3}
  The orthogonal group $O(2m+1)$ is the (transitive) group of isometries of
  $S^n$. Let $\Gamma$ be a subgroup that acts in a totally discontinuous manner on
  $S^n$. If $\gamma\in\Gamma$ has determinant $+1$, there exists an eigenvalue of
  $\gamma$ equal to $1$; then $\gamma$ has a fixed point, which implies $\gamma=e$.
  If $\tilde{\gamma}\in\Gamma$ has determinant $-1$, then $\tilde{\gamma}^2$ has
  determinant $1$, therefore $\tilde{\gamma}^2=e$. The eigenvalues of
  $\tilde{\gamma}$ are $1$ or $-1$. If $1$ is an eigenvalue then $\tilde{\gamma}=e$,
  contradicting $\det\tilde{\gamma}=-1$. As a result $\tilde{\gamma}=-e$, hence
  $\Gamma=\{e,-e\}$. Using \cref{prop:dc-ch8-4-3} we obtain that $M^n$ is either $S^n$, if
  $\Gamma=\{e\}$, or $P^n$, if $\Gamma=\{e,-e\}$. $\square$
\end{proof}

\begin{proposition}[metric of constant negative curvature on surfaces of genus $>1$]
  \label{prop:dc-ch8-4-5}
  \dcref{ch8:4.5}
  \notready
  Every compact orientable surface of genus $p>1$ can be provided with a metric of
  constant negative curvature.
\end{proposition}
\begin{proof}
  In the hyperbolic plane $H^2$ take a closed geodesic polygon $P$ with $4p$
  sides of equal lengths. By a well known process of identifying the sides, $P$ can
  be identified with a topological surface $M^2$ (Massey, Chap. 1). Let $\Gamma$ be
  the subgroup of isometries of $H^2$ generated by the isometries that identify the
  sides of $P$. The transforms of $P$ by $\Gamma$ "tile $H^2$ without any gaps" iff
  the sum $\alpha$ of the interior angles at the vertices of $P$ equals $2\pi$; in
  this case $M^2=H^2/\Gamma$ has a metric of constant curvature $-1$. We assert that
  it is possible to find a polygon $P$ satisfying $\alpha=2\pi$. By the Gauss–Bonnet
  Theorem, $\alpha=-A+2\pi(2p-1)$, where $A$ is the area of $P$ in the hyperbolic
  metric. If $P$ is arbitrarily small, $\alpha$ is arbitrarily close to $2\pi(2p-1)$;
  if we enlarge the area $A$, $\alpha$ becomes arbitrarily small. By continuously
  deforming $P$, there exists a geodesic polygon such that $\alpha=2\pi$. $\square$

  The spaces of constant curvature have an important role in the historical
  development of Riemannian Geometry, due to their relationship with non-Euclidean
  geometry. A non-Euclidean geometry is a complete Riemannian manifold $M$ together
  with a transitive group of isometries $G$ (the non-Euclidean motions) satisfying
  the \emph{Axiom of free mobility}: if $p,\tilde p\in M$, if $\gamma_1,\gamma_2$ are
  geodesics starting at $p$ and forming an angle $\alpha$, and if
  $\tilde\gamma_1,\tilde\gamma_2$ are geodesics starting at $\tilde p$ and forming
  the same angle $\alpha$, then there is a $g\in G$ with $g(p)=\tilde p$,
  $g(\gamma_1)=\tilde\gamma_1$, and $g(\gamma_2)=\tilde\gamma_2$. This corresponds
  to the side-angle-side condition for congruence of triangles in Euclidean geometry
  and implies that $M$ has constant sectional curvature, so the spaces of
  non-Euclidean Geometry are included among the space forms. The cases of $S^n$,
  $P^n$ and $H^n$ are called \emph{spherical}, \emph{elliptic} and \emph{hyperbolic} geometry,
  respectively.
\end{proof}

\section{Isometries of the hyperbolic space; Theorem of Liouville}

The isometries of hyperbolic space are closely related to the conformal
transformations of $\mathbb{R}^n$. A map $f:U\subset\mathbb{R}^n\to\mathbb{R}^n$ of
an open set is called \emph{conformal} if the (non-oriented) angles of intersecting
curves are preserved. The principal objective of this section is to show that the
isometries of the upper half-space $H^n\subset\mathbb{R}^n$ with the metric
$\delta_{ij}/x_n^2$ are the restrictions to $H^n$ of conformal transformations of
$\mathbb{R}^n$ that map $H^n$ onto $H^n$. For $\mathbb{R}^2\cong\mathbb{C}$, the
conformal transformations are the holomorphic or anti-holomorphic functions with
non-zero derivatives. For $\mathbb{R}^n$, $n>2$, the situation is radically
different, and being conformal imposes strong restrictions (Theorem of Liouville).

A necessary and sufficient condition for $f:U\subset\mathbb{R}^n\to\mathbb{R}^n$ to
be conformal is that for all $p\in U$ and all pairs of vectors $v_1,v_2$ at $p$:

$$
\langle df_p(v_1),df_p(v_2)\rangle=\lambda^2(p)\langle v_1,v_2\rangle,\qquad\lambda^2\neq 0.\qquad(3)
$$

Indeed, if $(3)$ holds then $|df_p(v)|^2=\lambda^2|v|^2$, hence

$$
\cos\angle(df_p(v_1),df_p(v_2))=\cos\angle(v_1,v_2),\qquad(4)
$$

so angles are preserved; conversely $(4)$ implies $(3)$. The positive function
$\lambda:U\to\mathbb{R}$ defined in $(3)$ is called the *coefficient of
conformality* of $f$.

\begin{definition}[conformal map with a coefficient of conformality]
  \label{def:dc-ch8-5-conformal}
  \lean{Riemannian.IsConformalWithCoeff,
    Riemannian.IsConformalWithCoeff.isConformalMap,
    Riemannian.IsConformalWithCoeff.conformalAt}
  \leanok
  A map $f:U\subset\mathbb{R}^n\to\mathbb{R}^n$ is \emph{conformal at $p$ with
  coefficient of conformality $\lambda(p)>0$} if $f$ is differentiable at $p$ and
  its differential rescales every inner product by $\lambda^2$, i.e. condition
  $(3)$ holds:
  $\langle df_p(v_1),df_p(v_2)\rangle=\lambda^2(p)\langle v_1,v_2\rangle$ for all
  $v_1,v_2$. Equivalently $df_p$ is a conformal linear map, so $f$ is conformal at
  $p$ in the sense of preserving non-oriented angles.
\end{definition}

\begin{lemma}[a conformal map preserves norms up to $\lambda$ and preserves angles: $(3)\Rightarrow(4)$]
  \label{lem:dc-ch8-5-conformal-angle}
  \uses{def:dc-ch8-5-conformal}
  \lean{Riemannian.IsConformalWithCoeff.norm_fderiv_eq,
    Riemannian.IsConformalWithCoeff.angle_fderiv_eq}
  \leanok
  If $f$ is conformal at $p$ with coefficient $\lambda$, then $|df_p(v)|=\lambda|v|$
  for every $v$, and $df_p$ preserves the non-oriented angle between any two
  vectors: $\angle(df_p(v_1),df_p(v_2))=\angle(v_1,v_2)$, which is do Carmo's
  equation $(4)$.
\end{lemma}
\begin{proof}
  \leanok
  From $(3)$ with $v_1=v_2=v$, $|df_p(v)|^2=\langle df_p(v),df_p(v)\rangle
  =\lambda^2\langle v,v\rangle=\lambda^2|v|^2$, so $|df_p(v)|=\lambda|v|$ since
  $\lambda>0$. Hence the cosine of the angle
  $\langle df_p(v_1),df_p(v_2)\rangle/(|df_p(v_1)|\,|df_p(v_2)|)
  =\lambda^2\langle v_1,v_2\rangle/(\lambda^2|v_1|\,|v_2|)
  =\langle v_1,v_2\rangle/(|v_1|\,|v_2|)$
  is unchanged, giving $(4)$.
\end{proof}

\begin{example}[isometries, dilatations, inversions]
  \label{ex:dc-ch8-5-1}
  \dcref{ch8:5.1}
  \uses{def:dc-ch8-5-conformal}
  \lean{Riemannian.isConformalWithCoeff_isometry,
    Riemannian.isConformalWithCoeff_dilatation,
    Riemannian.isConformalWithCoeff_inversion}
  \leanok
  An isometry of $\mathbb{R}^n$ (orthogonal linear transformation followed by a
  translation) is a conformal transformation with coefficient of conformality
  $\lambda\equiv 1$. The linear transformation $f(p)=\lambda I(p)$, where $I$ is the
  identity matrix and $\lambda=\text{const.}>0$, is conformal with coefficient
  $\lambda$; $f$ is called a \emph{dilatation}. The \emph{inversion} with respect to the unit
  sphere centered at $p_0\in\mathbb{R}^n$, defined by

  $$
  f(p)=\frac{p-p_0}{|p-p_0|^2}+p_0,\qquad p\in\mathbb{R}^n-\{p_0\},\qquad(5)
  $$

  is a conformal transformation. Geometrically, $f$ takes $p$ to a point on the line
  through $p$ and $p_0$ at distance $1/|p-p_0|$ from $p_0$; $f$ keeps fixed the unit
  sphere around $p_0$ and permutes its interior with its exterior; in particular
  $f^2=\text{identity}$. For a vector $v$ at $p$,

  $$
  df_p(v)=\frac{v|p-p_0|^2-2\langle v,p-p_0\rangle(p-p_0)}{|p-p_0|^4},
  $$

  and one computes $|df_p(v)|^2=\langle v,v\rangle/|p-p_0|^4$, so the inversion $(5)$
  is conformal with coefficient of conformality $\lambda=1/|p-p_0|^2$.

  The Theorem of Liouville asserts the surprising fact that every conformal
  transformation $f$ of an open set $U\subset\mathbb{R}^n$, $n\geq 3$, extends to a
  composition of isometries, dilatations and inversions, at most one of each.
\end{example}

\begin{lemma}[conformal maps are closed under composition]
  \label{lem:dc-ch8-5-conformal-comp}
  \uses{def:dc-ch8-5-conformal}
  \lean{Riemannian.IsConformalWithCoeff.comp,
    Riemannian.isConformalWithCoeff_id}
  \leanok
  If $f$ is conformal at $p$ with coefficient $\lambda$ and $g$ is conformal at
  $f(p)$ with coefficient $\mu$, then $g\circ f$ is conformal at $p$ with
  coefficient $\mu\lambda$. In particular the identity is conformal with
  coefficient $1$, and any composition of isometries, dilatations and inversions —
  the class of maps in the Theorem of Liouville — is conformal.
\end{lemma}
\begin{proof}
  \leanok
  By the chain rule $d(g\circ f)_p=dg_{f(p)}\circ df_p$. For all $v_1,v_2$,
  $\langle d(g\circ f)_p(v_1),d(g\circ f)_p(v_2)\rangle
  =\mu^2\langle df_p(v_1),df_p(v_2)\rangle
  =\mu^2\lambda^2\langle v_1,v_2\rangle=(\mu\lambda)^2\langle v_1,v_2\rangle$,
  and $\mu\lambda>0$.
\end{proof}

\begin{theorem}[Liouville]
  \label{thm:dc-ch8-5-2}
  \dcref{ch8:5.2}
  \notready
  Let $f:U\to\mathbb{R}^n$, $n\geq 3$, be a conformal transformation of an open set
  $U\subset\mathbb{R}^n$. Then $f$ is the restriction to $U$ of a composition of
  isometries, dilatations or inversions, at most one of each.
\end{theorem}
\begin{proof}
  Let $a_1=(1,0,\dots,0),\dots,a_n=(0,\dots,0,1)$ be the canonical basis and
  $(x_1,\dots,x_n)$ the cartesian coordinates. Let $e_1,\dots,e_n$ be parallel
  differentiable vector fields on $U$ with $\langle e_i,e_j\rangle=\delta_{ij}$. If
  $\lambda$ is the coefficient of conformality of $f$,

  $$
  \langle df(e_i),df(e_k)\rangle=\lambda^2\delta_{ik},\qquad i,k=1,\dots,n.\qquad(6)
  $$

  Let $d^2 f$ be the second differential of $f$, a symmetric bilinear map with
  $d^2 f(a_i,a_j)=\partial^2 f/\partial x_i\partial x_j$. Differentiating $(6)$ with
  $i,j,k$ distinct and combining the three cyclic equations yields
  $\langle d^2 f(e_k,e_j),df(e_i)\rangle=0$ for $i,j,k$ distinct. Letting $i$ vary,
  $d^2 f(e_k,e_j)$ belongs to the plane generated by $df(e_j)$ and $df(e_k)$, so
  $d^2 f(e_k,e_j)=\mu\,df(e_k)+\nu\,df(e_j)$ with $\mu=d\lambda(e_j)/\lambda$,
  $\nu=d\lambda(e_k)/\lambda$, that is,

  $$
  d^2 f(e_k,e_j)=\frac{1}{\lambda}(df(e_k)\,d\lambda(e_j)+df(e_j)\,d\lambda(e_k)).\qquad(7)
  $$

  Put $\rho=1/\lambda$. Using $d(\rho f)=d\rho\,f+\rho\,df$ and $(7)$,

  $$
  d^2(\rho f)(e_k,e_j)=d^2\rho(e_k,e_j)f.\qquad(8)
  $$

  We claim $d^2\rho(e_k,e_j)=0$ for $k\neq j$. Computing the third differential
  $d^3(\rho f)$ via $(8)$ and using the symmetry of the first member in $i,j,k$, one
  obtains $d^2\rho(e_k,e_j)df(e_i)=d^2\rho(e_k,e_i)df(e_j)$; since $df(e_i)$ and
  $df(e_j)$ are linearly independent and $i,j,k$ are distinct, $d^2\rho(e_k,e_j)=0$
  for all $j\neq k$. Choosing $e_1,\dots,e_n$ to form a prescribed orthonormal basis
  at any $p$, this holds for every orthonormal basis; since
  $0=d^2\rho\big(\tfrac{e_j+e_k}{\sqrt 2},\tfrac{e_j-e_k}{\sqrt 2}\big)=\tfrac12\{d^2\rho(e_j,e_j)-d^2\rho(e_k,e_k)\}$,
  we get $d^2\rho(e_j,e_j)=d^2\rho(e_k,e_k)$ for $j\neq k$. In summary, for every
  $p\in U$ and any orthonormal basis, $d^2\rho(e_j,e_k)=\sigma\delta_{jk}$; in the
  canonical basis,

  $$
  \frac{\partial^2\rho}{\partial x_i\partial x_j}=\sigma\delta_{ij}.\qquad(9)
  $$

  Differentiating $(9)$ gives $\partial\sigma/\partial x_i=0$ ($i\neq j$), so
  $\sigma=\text{const.}$

  *Case $\sigma=\text{const.}\neq 0$.* Equation $(9)$ implies

  $$
  \rho=\frac{\sigma}{2}\sum x_i^2+\sigma\sum b_i x_i+c,\quad b_i,c\text{ constants},\qquad(10)
  $$

  which we write as $1/\lambda=\rho=a_1|p-p_0|^2+k_1$ with $a_1=\sigma/2$,
  $k_1=\text{const.}$, $p_0\in\mathbb{R}^n$. The proof in this case will be complete
  once we show $k_1=0$: if $k_1=0$, taking the inversion
  $g(p)=\frac{p-p_0}{|p-p_0|^2}+p_0$ and $h=g\circ f^{-1}$, $h$ is conformal with
  coefficient $a_1$, hence an isometry followed by a dilatation; thus
  $f=h^{-1}\circ g$ is an inversion followed by a dilatation, followed by an
  isometry. To show $k_1=0$, apply the argument to $f^{-1}$:
  $\lambda=a_2|f(p)-q_0|^2+k_2$, hence

  $$
  (a_1|p-p_0|^2+k_1)(a_2|f(p)-q_0|^2+k_2)=1.\qquad(11)
  $$

  Equation $(11)$ shows that a sphere of center $p_0$ is mapped by $f$ into a sphere
  of center $q_0$, and since $f$ preserves angles, radii map to radii. For a radial
  segment $p(s)$, $0\leq s\leq s_0$, of arc length $s$,

  $$
  \int_0^{s_0}\Big|df\Big(\frac{dp}{ds}\Big)\Big|\,ds=\int_0^{s_0}\frac{ds}{a_1|p(s)-p_0|^2+k_1}=|f(p(s_0))-f(p(0))|.
  $$

  If $k_1\neq 0$, the left side is a transcendental function of $|p(s_0)-p_0|$, while
  $(11)$ implies it is algebraic. This contradiction shows $k_1=0$, completing the
  case $\sigma\neq0$.

  *Case $\sigma=0$.* Here $\rho=1/\lambda=\sum a_i x_i+c_1=A_1(x)+c_1$. Applying the
  initial argument to $f^{-1}$,

  $$
  (A_1(x)+c_1)(A_2(f(x))+c_2)=1.\qquad(11')
  $$

  Equation $(11')$ shows a hyperplane parallel to $A_1=0$ maps into a hyperplane
  parallel to $A_2=0$; a line perpendicular to $A_1=0$ maps into a line. For a
  segment $p(s)$ of such a line,

  $$
  |f(p(s_0))-f(p(0))|=\int_0^{s_0}\frac{ds}{A_1(p(s))+c_1},
  $$

  which contradicts $(11')$ unless $A_1(p(s))\equiv 0$. We conclude that if
  $\sigma=0$, $\lambda=\text{const.}$; then lengths of tangent vectors are multiplied
  by a constant $\lambda$ and $f$ is an isometry followed by a dilatation. $\square$
\end{proof}

\begin{theorem}[isometries of $H^n$ are conformal transformations preserving $H^n$]
  \label{thm:dc-ch8-5-3}
  \dcref{ch8:5.3}
  \uses{def:dc-ch8-5-conformal,def:dc-ch8-3-hyperbolic-space}
  \notready
  The isometries of $H^n$ are the restrictions to $H^n\subset\mathbb{R}^n$ of the
  conformal transformations of $\mathbb{R}^n$ that take $H^n$ onto itself.
\end{theorem}
\begin{proof}
  \uses{cor:dc-ch8-2-3,thm:dc-ch8-5-2}
  Suppose first $n\geq 3$. Let $f:H^n\to H^n$ be an isometry in the metric
  $g_{ij}$. By Liouville's Theorem, $f$ extends to a conformal map of $\mathbb{R}^n$
  with the metric $\delta_{ij}$. Since $H^n$ is complete in $g_{ij}$, $f$ maps $H^n$
  onto $H^n$. Conversely, let $f:H^n\to H^n$ be a conformal transformation of $H^n$
  onto $H^n$ and let $e_1,\dots,e_n$ be an orthonormal basis (in $g_{ij}$) at
  $p\in H^n$. Since $g_{ij}=\delta_{ij}/x_n^2$ and $f$ is conformal, there exists
  $\lambda^2>0$ with $\langle df_p(e_i),df_p(e_j)\rangle=\lambda^2\delta_{ij}$, so
  $\{df_p(e_i)/\lambda\}$ is orthonormal at $f(p)$. \cref{cor:dc-ch8-2-3} of Cartan's
  Theorem, together with the global construction after Theorem 2.1 (the exponential
  maps of $H^n$ are global diffeomorphisms), gives an isometry $g$ of $H^n$ taking
  $p$ to $f(p)$ with $dg(e_i)=df(e_i)/\lambda$. By the first part, $g$ is conformal. Hence
  $h=g^{-1}\circ f$ is the restriction to $H^n$ of a conformal map of $\mathbb{R}^n$
  taking $H^n$ onto $H^n$, leaving $p$ fixed and satisfying $dh_p=\lambda I$. It
  remains to show $h=\text{identity}$.

  Let $P$ be a hyperplane through $p$. By Liouville's theorem, $h(P)$ is a hyperplane
  or a sphere through $p$; since $dh_p$ is a multiple of the identity, $P$ and $h(P)$
  are tangent at $p$. Since $h$ takes $\partial H^n$ into itself and is conformal,
  the angle of $P$ with $\partial H^n$ equals the angle of $h(P)$ with $\partial H^n$.
  We claim $h(P)=P$. Consider a line $r_1$ through $p$ perpendicular to
  $\partial H^n$, $q_1=r_1\cap\partial H^n$. Since $h(r_1)$ is a circle or a line
  making the same angle with $\partial H^n$, $h(r_1)=r_1$, hence $h(q_1)=q_1$. So if
  $P$ is perpendicular to $\partial H^n$, then $h(P)=P$. If $P$ is not perpendicular,
  let $r$ be a line in $P$; for $h(r)$ to be a circle making the same angle
  $\alpha$ it is necessary that $q_1\in h(r)$ (Fig. 2), contradicting $h(q_1)=q_1$;
  hence $h(r)$ is a line, $h(r)=r$, and $h(P)=P$. It follows that $h$ cannot be an
  inversion (the image of some plane would be a sphere) or an isometry distinct from
  the identity. From Liouville's theorem, $h$ is a dilatation; since $h$ takes $H^n$
  onto $H^n$, $h$ is the identity.

  For $n=2$, a simple calculation shows the conformal transformations of the form

  $$
  f(z)=\frac{az+b}{cz+d},\quad z\in H^2\subset\mathbb{C},\ a,b,c,d\in\mathbb{R},\ ad-bc>0\qquad(12)
  $$

  (which map $H^2$ onto $H^2$) are isometries of $H^2$ with the metric $g_{ij}$.
  Moreover, for a fixed point $p_0\in H^2$ and unit vector $v_0$ at $p_0$, there
  exists a transformation $(12)$ taking an arbitrary $(p,v)$ to $(p_0,v_0)$ ($p$ and
  $v$ are determined by three parameters, the number of parameters of $f$). Since
  there is a unique orientation-preserving isometry of $H^2$ taking $(p,v)$ to
  $(p_0,v_0)$, all orientation-preserving isometries of $H^2$ are of the form (12).
  Composing these with $z\mapsto-\bar z$ gives the orientation-reversing
  anti-holomorphic isometries, so together they give all isometries of $H^2$.
  $\square$

  To conclude this section, we identify some important hypersurfaces of $H^n$. The
  intersection with $H^n$ of hyperplanes of $\mathbb{R}^n$ orthogonal to $\partial H^n$,
  and of spheres with center on $\partial H^n$, are totally geodesic submanifolds of
  $H^n$ (Exercise 2). To determine intersections of planes and spheres in any
  position, it is convenient to use the \emph{ball model}: let
  $B^n=\{p\in\mathbb{R}^n;\ |p|<2\}$ with the metric

  $$
  h_{ij}(p)=\frac{\delta_{ij}}{(1-\tfrac14|p|^2)^2}.
  $$
\end{proof}

\begin{remark}[the ball model; horospheres and hyperspheres]
  \label{rem:dc-ch8-5-4}
  \dcref{ch8:5.4}
  \uses{def:dc-ch8-5-conformal,def:dc-ch8-3-hyperbolic-space,thm:dc-ch8-5-3}
  \notready
  $B^n$ is isometric to $H^n$ via the map

  $$
  f(p)=4\frac{p-p_0}{|p-p_0|^2}-(0,\dots,1),\qquad p_0=(0,\dots,-2),\qquad(13)
  $$

  since $\langle df_p(v),df_p(v)\rangle=16\langle v,v\rangle/|p-p_0|^4$ and
  $f_n(p)=(4-|p|^2)/|p-p_0|^2$ give
  $\langle df_p(v),df_p(v)\rangle/(f_n(p))^2=\langle v,v\rangle/(1-\tfrac14|p|^2)^2$.
  Because $f$ is injective, $f$ is an isometry of $B^n$ into $H^n$, taking
  $\partial B^n-\{p_0\}$ into $\partial H^n$. Conjugating by (13) and applying
  \cref{thm:dc-ch8-5-3}, $g:B^n\to B^n$ is an isometry in $h_{ij}$ iff $g$ is the
  restriction of a conformal transformation of $\mathbb{R}^n$ taking $B^n$ onto
  $B^n$.

  Let $S\subset H^n$ be an $(n-1)$-sphere of Euclidean space completely contained in
  $H^n$. Then $S$ is a \emph{geodesic sphere} of $H^n$: its image $f^{-1}(S)\subset B^n$ is
  a Euclidean sphere, which can be mapped onto a sphere centered at the origin by an
  isometry of $B^n$; since $h_{ij}$ is symmetric with respect to the origin, the
  result is a geodesic sphere, and so is $S$.

  Let $S$ be an $(n-1)$-sphere tangent to $\partial H^n$ at $p$ with
  $S-\{p\}\subset H^n$. By an inversion of $\mathbb{R}^n$ at $p$ (an isometry of
  $H^n$), $S$ is mapped to a hyperplane $P$ parallel to $\partial H^n$. The induced
  metric on $P$ is a multiple of the Euclidean metric, so $P$ has constant curvature
  zero, as does $S-\{p\}$. Such submanifolds are called \emph{horospheres} of $H^n$.

  Let $S$ be a Euclidean sphere cutting $\partial H^n$ at an angle $\alpha$, and
  denote $S\cap H^n=\Sigma$. Through an inversion at a point of $S\cap\partial H^n$,
  $\Sigma$ is mapped isometrically into the intersection with $H^n$ of a hyperplane $P$
  cutting $\partial H^n$ at the same angle $\alpha$. Let $Q$ be the hyperplane
  orthogonal to $\partial H^n$ containing $P\cap\partial H^n$. Then $P$ is a
  hypersurface equidistant from the totally geodesic hypersurface $Q$: for a geodesic
  $\gamma_r$ (a semi-circle of radius $r$ with center $0\in P\cap\partial H^n$,
  lying in the plane perpendicular to $P\cap\partial H^n$), a
  homothety with center $0$ takes the circle of radius $r$ into one of any radius, so
  the length of $\gamma_r$ between its intersections with $P$ and $Q$ does not depend
  on $r$. We conclude that $P$, or its isometric image $\Sigma$, is obtained by taking
  geodesics perpendicular to $Q$ and marking a fixed distance. Such hypersurfaces are
  called \emph{equidistant surfaces} (or \emph{hyperspheres}). In Exercise 6 it is shown that
  the hypersurfaces above are exactly the umbilic ones, and their mean and sectional
  curvatures are calculated.
\end{remark}
